\documentclass[reqno]{amsart}
\usepackage{a4wide}
\usepackage{hyperref}

\AtBeginDocument{{\noindent\small
\emph{Electronic Journal of Differential Equations},
Vol. 2026 (2026), No. 06, pp. 1--20.\newline
ISSN: 1072-6691. URL: https://ejde.math.txstate.edu, DOI: 10.58997/ejde.2026.06}
\thanks{\copyright 2026. This work is licensed under a CC BY 4.0 license.}
\vspace{8mm}}

\begin{document}
\title[\hfilneg EJDE-2026/06\hfil 
Kirchhoff type system with van der Waals type potentials]
{Multiple positive normalized solutions for Kirchhoff type system with
van der Waals type potentials}

\author[Z. Chen, M. Li \hfil EJDE-2026/06\hfilneg]
{Zhewen Chen, Muzi Li}

\address{Zhewen Chen (corresponding author) \newline
College of Science,
Jimei University,
Xiamen 361021, China}
\email{zhewenchen@jmu.edu.cn}

\address{Muzi Li \newline
College of Mathematics and Statistics,
Fujian Normal University,
Qishan Campus, Fuzhou 350117, China}
\email{mutoudededemuzi@163.com}

\thanks{Submitted July 2, 2025. Published January 19, 2026.}
\subjclass[2020]{35J20, 35J47, 35J50}
\keywords{Ground state; normalized solution; Van der Waals type potential;
Kirchhoff system}

\begin{abstract}
This article shows the existence of normalized solutions for Kirchhoff type
system with van der Waals type potentials,
\begin{gather*}
-(a+b\int_{\mathbb{R}^N}|\nabla u_1|^2dx)\Delta u_1=\lambda_1 u_1+\mu_1(I_\alpha\ast|u_1|^{p_1})|u_1|^{p_1-2}u_1+
\Theta r_1(I_\beta\ast|u_2|^{r_2})|u_1|^{r_1-2}u_1, \\
-(a+b\int_{\mathbb{R}^N}|\nabla u_2|^2dx)\Delta u_2=\lambda_2 u_2+\mu_2(I_\alpha\ast|u_2|^{p_2})|u_2|^{p_2-2}u_2+
\Theta r_2(I_\beta\ast|u_1|^{r_1})|u_2|^{r_2-2}u_2, \\
\int_{\mathbb{R}^N}|u_1|^2dx=d_1>0,\quad
\int_{\mathbb{R}^N}|u_2|^2dx=d_2>0,
\end{gather*}
where $N=3,4$, $\mu_1,\mu_2,\Theta>0$, $\frac{N+\alpha}{N}<p_1,p_2<\frac{N+\alpha+2}{N}$,
$2\cdot\frac{N+\beta}{N}<r_1+r_2<2\cdot2_\beta^*=2\cdot\frac{N+\beta}{N-2}$,
$0<\alpha,\beta<N$, $I_\alpha$ and $I_\beta$ are the Riesz potentials.
We show that the system has a positive least energy solution at negative energy
level for $\Theta$ small. In addition, we also prove that the system admits a
high energy positive solution at positive energy level in the special case.
\end{abstract}

\maketitle
\numberwithin{equation}{section}
\newtheorem{theorem}{Theorem}[section]
\newtheorem{lemma}[theorem]{Lemma}
\newtheorem{definition}[theorem]{Definition}
\newtheorem{remark}[theorem]{Remark}
\allowdisplaybreaks

\section{Introduction and main results}

In this article, we study the existence of the positive normalized solutions for
Kirchhoff type system with van der Waals type potentials
\begin{equation}\label{yy}
\begin{gathered}
-(a+b\int_{\mathbb{R}^N}|\nabla u_1|^2dx)\Delta u_1
=\lambda_1 u_1+\mu_1(I_\alpha\ast|u_1|^{p_1})|u_1|^{p_1-2}u_1+
 \Theta r_1(I_\beta\ast|u_2|^{r_2})|u_1|^{r_1-2}u_1, \\
-(a+b\int_{\mathbb{R}^N}|\nabla u_2|^2dx)\Delta u_2
=\lambda_2 u_2+\mu_2(I_\alpha\ast|u_2|^{p_2})|u_2|^{p_2-2}u_2+
 \Theta r_2(I_\beta\ast|u_1|^{r_1})|u_2|^{r_2-2}u_2,
\end{gathered}
\end{equation}
with the $L^2$-mass constraint
\begin{equation}\label{yues}
 \int_{\mathbb{R}^N}|u_1|^2dx=d_1>0,\quad
 \int_{\mathbb{R}^N}|u_2|^2dx=d_2>0,
\end{equation}
where $N=3,4$, $a,b>0$, $\alpha,\beta\in(0,N)$, $I_\alpha,I_\beta$ are the Riesz
potentials defined for every $x\in\mathbb{R}^N\backslash\{0\}$ by
\begin{equation}\label{gamma}
 I_\alpha(x):=\frac{A_\alpha(N)}{|x|^{N-\alpha}},\quad
 A_\alpha(N):=\frac{\Gamma(\frac{N-\alpha}{2})}{\Gamma(\frac{\alpha}{2})\pi^{N/2}
 2^\alpha}
\end{equation}
with $\Gamma$ denoting the Gamma function. Throughout this paper, we always require
that $a,b,\mu_1,\mu_2,\Theta>0$, and assume that $2<p_1,p_2
<2+\frac{4}{N}$, $r_1,r_2>\frac{N+\beta}{N}$,
$2\cdot\frac{N+\beta}{N}<r_1+r_2<2\cdot2_\beta^*$.

Because of the appearance of the term $\int_{\mathbb{R}^N}|\nabla u|^2dx$,
\eqref{yy} is regard as a nonlocal problem, which implies that equation \eqref{yy}
is not a pointwise identity. Moreover, this phenomenon also leads to some mathematical
difficulties that make the study of \eqref{yy} more interesting.
Problem \eqref{yy} originates
from the stationary analog of the equation
\[
\rho\frac{\partial^2 u}{\partial t^2}
-\Big(\frac{Y_0}{k}+\frac{R}{2J}\int_{0}^{Q}|\frac{\partial u}{\partial x}|^2dx\Big)
\frac{\partial^2 u}{\partial x^2}=0,
\]
which was proposed by Kirchhoff \cite{Kf4} in 1883, and being an extension of the
classical D'Alembert's wave equations for free vibration of elastic strings.
Kirchhoff's model takes into account the changes in length of the string produced by
transverse vibrations. In the last decades, Because of the strong background in physics,
starting with the framework of Lions \cite{Li4}, many mathematicians have established
many interesting conclusions about Kirchhoff-type problems
\cite{Kf1,Kf2,Kf3,Kf4,Kf5,Kf6,ygg,guo1,guo2}.

In \eqref{yy}, if $\lambda\in\mathbb{R}$ is fixed, then we call \eqref{yy} the fixed
frequency problem. One can adopt the traditional variational method,
looking for critical points of $F_\lambda(u_1,u_2)$, or fixed point theory,
bifurcation, topological methods, Nehari manifold method and Lyapunov-Schmidt reduction,
where
\begin{align*}
F_\lambda(u_1,u_2)
&:=\frac{a}{2}\sum_{i=1}^2\| \nabla u_i\|_2^2
 +\frac{b}{4}\sum_{i=1}^2\| \nabla u_i\|_2^4-\sum_{i=1}^2\lambda_i\| u_i\|_2^2
 -\sum_{i=1}^2\frac{\mu_i}{2p_i}\int_{\mathbb{R}^N}
 (I_\alpha\ast|u_i|^{p_i})|u_i|^{p_i}dx\\
&\quad -\Theta\int_{\mathbb{R}^N}(I_\beta\ast|u_1|^{p_1})|u_2|^{p_2}dx.
\end{align*}
In recent decades, because of the application to physics, mathematicians are
interested in solutions that satisfy $L^2-$mass constraint \eqref{yues}.
In this direction, the mass $d_1,d_2>0$ is prescribed, the frequency $\lambda_i$
cannot be determined a priori, but is a part of unknown which appears as Lagrange
multipliers.
In this case, mathematicians often call \eqref{yy}-\eqref{yues} the fixed mass
problem and the solution is called a normalized solution.
One can get a normalized solution to problem \eqref{yy} by looking for a critical
point of the  functional
\begin{align*}
&F(u_1,u_2) \\
&:=\frac{a}{2}\sum_{i=1}^2\| \nabla u_i\|_2^2
 +\frac{b}{4}\sum_{i=1}^2\| \nabla u_i\|_2^4-
\sum_{i=1}^2\frac{\mu_i}{2p_i}\int_{\mathbb{R}^N}(I_\alpha\ast|u_i|^{p_i})|u_i|^{p_i}dx
-\Theta\int_{\mathbb{R}^N}(I_\beta\ast|u_1|^{p_1})|u_2|^{p_2}dx
\end{align*}
constrained on $S(d_1)\times S(d_2)$, where
$S(d):=\{u\in H^1(\mathbb{R}^N): \int_{\mathbb{R}^N}|u|^2dx=d>0\}$.
It is standard to check that $F\in C^1$.

For the fixed mass problem, the $L^2$-mass constraint \eqref{yues} presents some
mathematical difficulties. As opposed to the fixed
frequency problem, the fixed mass problem will have many technical difficulties when
dealing with it in a variational framework:
(I) the Nehari manifold method is inapplicable;
(II) the Lagrange multipliers must be controlled;
(III) for the fixed frequency problem, usually a
nontrivial weak limit is also a solution.
However, for the fixed mass problem, even though the weak limit is nontrivial, the
$L^2$-mass constraint may be not satisfied;
(IV) the $L^2$-mass critical exponent seriously affects the geometric structure of the
functional.

The equation
\begin{equation}\label{yyu5}
 i\partial_t\varPsi+\Delta \varPsi+(V\ast|\varPsi|^p)|\varPsi|^{p-2}\varPsi=0,    \text{in} \mathbb{R}^+\times\mathbb{R}^N,\\
 \end{equation}
has $\varPsi:\mathbb{R}^+\times\mathbb{R}^N\to \mathbb{C}$ is a complex
valued function, $V(x)=\delta_{\zeta}\frac{1}{|x|^\zeta}\pm
 \delta_{\theta}\frac{1}{|x|^\theta}$($\alpha,\beta>1$)
is the van der Waals type potential (see \cite{VWp3,BZd27}).
The van der Waals coefficients
$\delta_6$, $\delta_8$ and $\delta_{10}$ of alkaline-earth interactions
calculated by Porsev and Derevianko using relativistic many-body perturbation
theory are believed to be accurate to $1/100$ (see\cite{cfrv}).
As one of the van der Waals type potentials the Lennard-Jones potential
\[
 V_{LJ}(\nu)=\pm\xi(\frac{1}{\nu^\zeta }-\frac{1}{\nu^\theta})
\]
with $\zeta =6$ and $\theta=12$, is often used as an approximate model for the
isotropic part of a total van der
 Waals force as a function of distance (see \cite{VWp3,BZd27}).

When $a=1$, $b=0$ and the response function is a delta function, i.e.\
$I_\alpha(x)=I_\beta(x)=\delta(x)$, the nonlinear response is local and
problem \eqref{yy} with
prescribed mass turns out to be
\begin{equation}\label{hend}
\begin{gathered}
-\Delta u_1=\lambda_1 u_1+\mu_1|u_1|^{p_1-2}u_1+
 \Theta r_1|u_2|^{r_2}|u_1|^{r_1-2}u_1, \\
-\Delta u_2=\lambda_2 u_2+\mu_2|u_2|^{p_2-2}u_2+
 \Theta r_2|u_1|^{r_1}|u_2|^{r_2-2}u_2, \\
\end{gathered}
\end{equation}
with the $L^2$-mass constraint
$\int_{\mathbb{R}^N}|u_1|^2dx=d_1>0$, $\int_{\mathbb{R}^N}|u_2|^2dx=d_2>0$.
In recent years, mathematicians have drawn rich conclusions
concerning the existence, multiplicity and qualitative properties of
the normalized solutions of system \eqref{hend}
(see \cite{Huij1,Huij2,Huij3,Huij4,Huij5,Huij6}).
Gou et al.\ \cite{Huij6} showed that
the system admits two normalized solutions under the conditions that
$2<p_1,p_2<2+\frac{4}{N}<r_1+r_2<2^*=\frac{2N}{N-2}$
(resp.\ $2+\frac{4}{N}<r_1+r_2<2+\frac{4}{N}<p_1,p_2<2^*$)
by using minimizing methods, Poho\v{z}aev type manifold, Schwarz rearrangements
and mountain pass theorem.
Bartsch et al.\ \cite{Huij4} focused on the choice $p_1=p_2=2^*$, and allowing
$r_1+r_2$ to be mass-subcritical, mass-critical or mass-supercritical.
The authors proved the existence and non-existence of normalized ground state for
different ranges of $\Theta$.

When $a,b>0$ and the response function is a delta function,
i.e.\ $I_\alpha(x)=I_\beta(x)=\delta(x)$,
system \eqref{yy} becomes
\begin{equation}\label{sss8}
\begin{gathered}
 -(a+b\int_{\mathbb{R}^N}|\nabla u_1|^2dx)\Delta u_1=\lambda_1 u_1+\mu_1|u_1|^{p_1-2}u_1+
 \Theta r_1|u_2|^{r_2}|u_1|^{r_1-2}u_1   x\in\mathbb{R}^N,\\
-(a+b\int_{\mathbb{R}^N}|\nabla u_2|^2dx)\Delta u_2=\lambda_2 u_2+\mu_2|u_2|^{p_2-2}u_2+
 \Theta r_2|u_1|^{r_1}|u_2|^{r_2-2}u_2   x\in\mathbb{R}^N.\\
 \end{gathered}
\end{equation}
When $N\leq3$ and $2<p_1,p_2,r_1+r_2<2+\frac{8}{N}$, Cao et al.\
 \cite{kf2} proved that system \eqref{sss8} admits a positive
normalized solution. When $N=2,3$, Yang \cite{Kf3} showed that system
\eqref{sss8} admits a normalized ground state under the condition that
$2+\frac{8}{N}<p_1,p_2,r_1+r_2<2^*$ and $2<r_1+r_2<2+\frac{8}{N}<p_1,p_2<2^*$,
respectively.

When $a=1$, $b=0$ and the response function is a delta function, i.e.\
$I_\alpha(x)=I_\beta(x)=\delta(x)$,
system \eqref{yy} becomes the   Choquard system (Hartree system)
\begin{equation}\label{ChqH}
 \begin{gathered}
-\Delta u_1=\lambda_1 u_1+\mu_1(I_\alpha\ast|u_1|^{p_1})|u_1|^{p_1-2}u_1
+\Theta r_1(I_\alpha\ast|u_2|^{r_2})|u_1|^{r_1-2}u_1   x\in\mathbb{R}^N,\\
-\Delta u_2=\lambda_2 u_2+\mu_2(I_\alpha\ast|u_2|^{p_2})|u_2|^{p_2-2}u_2
+\Theta r_2(I_\alpha\ast|u_1|^{r_1})|u_2|^{r_2-2}u_2   x\in\mathbb{R}^N.
 \end{gathered}
\end{equation}
When $\mu_1,\mu_2,\Theta>0$ and $\frac{N+\alpha}{N}<r_1,r_2$, Geng et al.\
\cite{JFx} proved that system \eqref{ChqH}
has a normalized ground state under the condition that $N\geq3$ and
$\frac{N+\alpha}{N}<p_1,p_2<\frac{N+\alpha+2}{N}<r_1,r_2<2_\alpha^*$
by using Schwartz rearrangement. In addition, they proved that the system
\eqref{ChqH} has a second solution
for $N=3$ and $p_1=p_2=\alpha=2$.
They also proved that  \eqref{ChqH} has a second solution
for $N=5$, $\frac{N+\alpha}{N}<r_1,r_2<\frac{N+\alpha+2}{N}<p_1,p_2<2_\alpha^*$ and
$p_1=p_2=\alpha=2$. For
Hardy-Littlewood-Sobolev critical case, i.e.\ $p_1=p_2=2_\alpha^*$, Zhang et al.\
\cite{linjie} proved that system \eqref{ChqH}
has a normalized ground state for different ranges of $\Theta$ when $r_1+r_2$ is set
to be mass subcritical, mass critical and mass supercritical, respectively.

However, as far as we know, for the case of $a,b>0$ and the response functions are
different Riesz potential functions,
the existence of the solution of the system is still unknown.

\begin{definition}\label{df82} \rm
 $(u_1,u_2)$ is a \emph{normalized ground state} to \eqref{yy}-\eqref{yues}
if $F'|_{S(d_1)\times S(d_2)}(u_1,u_2)=0$ and
\[
 F(u_1,u_2)=\inf\{F(v_1,v_2): (v_1,v_2)\in S(d_1)\times S(d_2),
  F'|_{S(d_1)\times S(d_2)}(v_1,v_2)=0\}.
\]
Furthermore, $(w_1,w_2)$ is a \emph{high energy normalized solution}
to \eqref{yy}-\eqref{yues} if $F'|_{S(d_1)\times S(d_2)}(w_1,w_2)=0$ and
\[
 F(w_1,w_2)>\inf\{F(v_1,v_2): (v_1,v_2)\in S(d_1)\times S(d_2),
 F'|_{S(d_1)\times S(d_2)}(v_1,v_2)=0\}.
\]
\end{definition}

In this article, we demonstrate our main results in the following two cases:
\begin{itemize}
\item[(A1)] $\frac{N+\alpha}{N}<p_1,p_2<\frac{N+\alpha+2}{N}$,
$r_1,r_2>\frac{N+\beta}{N}$ and
$2\cdot\frac{N+\beta}{N}<r_1+r_2<2\cdot\frac{N+\beta+4}{N}$.

\item[(A2)]  $\frac{N+\alpha}{N}<p_1,p_2<\frac{N+\alpha+2}{N}$,
$r_1,r_2>\frac{N+\beta}{N}$ and $2\cdot\frac{N+\beta+4}{N}<r_1+r_2<2\cdot2_\beta^*$.
\end{itemize}
To restore some compactness, we search for critical points of $F$ constrained on
$S_r(d_1)\times S_r(d_2)$,
where $S_r(d):=\{u\in H_r^1(\mathbb{R}^N):\| u\|_2^2=d\}$ and
\[
H_r^1(\mathbb{R}^N)=\{u\in H^1(\mathbb{R}^N):
 u \text{is a radially symmetric function}.
\] 
For case of ($A_2$), $F$ is unbounded below in $S_r(d_1)\times S_r(d_2)$.
Hence, we need the Poho\v{z}aev manifold
$\mathcal{P}(d_1,d_2):=\{(u_1,u_2)\in S_r(d_1)\times S_r(d_2): P(u_1,u_2)=0\}$, where
\begin{align*}
P(u_1,u_2)
&=a\sum_{i=1}^2\| \nabla u_i\|_2^2+b\sum_{i=1}^2\| \nabla u_i\|_2^4
-\sum_{i=1}^2\frac{\mu_i}{2p_i}(Np_i-N-\alpha)
 \int_{\mathbb{R}^N}(I_\alpha\ast|u_i|^{p_i})|u_i|^{p_i}dx\\
&\quad-\Theta\frac{Nr-2(N+\beta)}{2}\int_{\mathbb{R}^N}(I_\beta\ast|u_1|^{r_1})
|u_2|^{r_2}dx
\end{align*}
with $r=r_1+r_2$. We define $W(k):=\{(u_1,u_2)\in S_r(d_1)\times
S_r(d_2):\| \nabla u_1\|_2^2+\| \nabla u_2\|_2^2<k\}$
and
\begin{equation}\label{jxh25}
 \sigma(d_1,d_2):=\inf_{W(k)}F(u_1,u_2)<0.
\end{equation}
For each $u\in S_r(d)$, we define
\[
 \tau\star u_i:=e^{\frac{N}{2}\tau}u_i(e^{\tau}x), \quad i=1,2.
\]
Hence,
\begin{equation}
\begin{aligned}
\Psi_{u_1,u_2}(\tau)
&:=F(\tau\star u_1,\tau\star u_2)\\
&=\frac{a}{2}e^{2\tau}\sum_{i=1}^2\|\nabla u_i\|_2^2
 +\frac{b}{4}e^{4\tau}\sum_{i=1}^2\|\nabla u_i\|_2^4\\
&\quad-\sum_{i=1}^2\frac{\mu_i}{2p_i}e^{(Np_i-N-\alpha)\tau}
 \int_{\mathbb{R}^N}(I_\alpha\ast|u_i|^{p_i})|u_i|^{p_i}dx\\
&\quad -\Theta e^{\frac{Nr}{2}-N-\beta}
 \int_{\mathbb{R}^N}(I_\beta\ast|u_1|^{r_1})|u_2|^{r_2}dx.
\end{aligned}
\end{equation}
By direct computations, we have $(\Psi_{u_1,u_2})'(0)=P(u_1,u_2)$.
We divide $\mathcal{P}(d_1,d_2)$ into 3 parts,
\[
\mathcal{P}(d_1,d_2)= \mathcal{P}_{-}(d_1,d_2)\cup\mathcal{P}_{0}(d_1,d_2)
\cup\mathcal{P}_{+}(d_1,d_2),
\]
 where
\begin{gather*}
\mathcal{P}_{-}(d_1,d_2)=\{(u_1,u_2)\in\mathcal{P}(d_1,d_2):(\Psi_{u_1,u_2})''(0)<0\},\\
\mathcal{P}_{0}(d_1,d_2)=\{(u_1,u_2)\in\mathcal{P}(d_1,d_2):(\Psi_{u_1,u_2})''(0)=0\},\\
\mathcal{P}_{+}(d_1,d_2)=\{(u_1,u_2)\in\mathcal{P}(d_1,d_2):(\Psi_{u_1,u_2})''(0)>0\},
\end{gather*}
and
\begin{align*}
(\Psi_{u_1,u_2})''(0)
&=2a\sum_{i=1}^2\|\nabla u_i\|_2^2+4b\sum_{i=1}^2\|\nabla u_i\|_2^4 \\
&\quad -\sum_{i=1}^2\frac{\mu_i}{2p_i}(Np_i-N-\alpha)^2
\int_{\mathbb{R}^N}(I_\alpha\ast|u_i|^{p_i})|u_i|^{p_i}dx \\
&\quad -\Theta(\frac{Nr}{2}-N-\beta)^2
 \int_{\mathbb{R}^N}(I_\beta\ast|u_1|^{r_1})|u_2|^{r_2}dx.
\end{align*}

\begin{theorem} \label{thm1.1}
Assume that {\rm (A1)} holds. Then, every minimizing sequence of \eqref{jxh25}
is compact, up to translation, in
$H^1_r(\mathbb{R}^N)\times H^1_r(\mathbb{R}^N)$.
Moreover, system \eqref{yy} has a positive normalized ground state.
\end{theorem}

\begin{theorem} \label{thm1.2}
Assume that {\rm (A2)} holds and $0<\alpha-2\leq\beta<\alpha<N$.
Then, there exist $k_0=k_0(d_1,d_2)>0$, $\Theta_*=\Theta_*(d_1,d_2)>0$, such that
for any $0<\Theta\leq\Theta_*$,
system \eqref{yy}-\eqref{yues} has a positive radial solution
$(\upsilon_1,\upsilon_2)\in W(k_0)$ at negative
level $F(\upsilon_1,\upsilon_2)<0$ for some $\lambda_1,\lambda_2<0$.
 \end{theorem}

\begin{theorem} \label{thm1.3}
Assume that {\rm (A2)} holds,and $a=1$, $b=0$, $N=3$ and
$p_1=p_2=\alpha=\beta=r_1=r_2=2$. Then, there exist
$k_0= k_0(d_1,d_2)>0$, $\Theta_*=\Theta_*(d_1,d_2)>0$, such that for any
$0<\Theta\leq\Theta_*$,
system \eqref{yy}-\eqref{yues} has a second positive radial solution
$(u_1,u_2)\in W(k_0)$ at positive level $F(u_1,u_2)>0$
 for some $\lambda_1,\lambda_2<0$.
\end{theorem}

% page 4

 \begin{remark}  \label{rmk1.1} \rm
Definition \ref{df82}, Theorem \ref{thm1.2} and \ref{Thm1.3}
indicate that \eqref{yy}-\eqref{yues} admit a normalized ground state at negative
level $F(\upsilon_1,\upsilon_2)<0$ and a high energy normalized solution at positive 
level $F(u_1,u_2)>0$, under the condition that
(A2), $a=1$, $b=0$, $N=3$ and $p_1=p_2=\alpha=\beta=r_1=r_2=2$ hold.
\end{remark}

\begin{remark} \label{rmk1.2} \rm
For nonlinear classical Choquard system \ref{ChqH}, the mass critical exponent is 
$\frac{N+\alpha+2}{N}$. In Kirchhoff-Choquard system,
by the appearance of the $\int_{\mathbb{R}^N}|\nabla u|^4dx$ term, the mass critical 
exponent becomes $\frac{N+\alpha+4}{N}$. In this
paper, if $\frac{N+\alpha}{N}<p_1,p_2<\frac{N+\alpha+4}{N}$, then Lemma 
\ref{ddh3} and Lemma \ref{zlx} do not hold.
\end{remark}

\begin{remark} \label{rmk1.3} \rm
Since we do not know whether the inequality
 \[
\| \nabla \{u_1,u_2\}^*\|_2^4\leq\| \nabla u_1\|_2^4+\| \nabla u_2\|_2^4
 \]
holds, we can not use the Schwartz rearrangement method to prove the compactness
of minimized sequence in this paper.
\end{remark}


From a physical point of view, it is of great importance to study the solution of 
 problem  \eqref{yy}, as confirmed in  \cite{VWp3,BZd27,cfrv}.
 We emphasize that this study seems to be the first contribution regarding existence 
of normalized ground states for a Kirchhoff system with van der Waals type potentials.

The rest of this article organized as follows: 
In Section 2, we present some preliminaries. 
In Section 3, we prove the existence
 of normalized ground states under the purely mass subcritical case. 
In Section 4, we prove the existence of the
 first solution, which is a local minimizer. 
In Section 5, the second solution is proved by using mountain pass theorem.
\smallskip

\noindent\textbf{Notation:} 
 $L^s(\mathbb{R}^N)$ is the Lebesgue space with the norms 
 $\| u\| _s=(\int_{\mathbb{R}^N}|u|^sdx)^{1/s}$,  $1<s<\infty$.\\
$H^1(\mathbb{R}^N)$ is the usual Sobolev space with norm 
$\| u\|_{H^1(\mathbb{R}^N)}=
 (\int_{\mathbb{R}^N}|\nabla u|^2+|u|^2dx)^{1/2}$.



\section{Preliminaries}

\begin{lemma}\label{L22}[Hardy-Littlewood-Sobolev inequality \cite{lieb}]
Let $N\geq1$, $p, r>1$, and $0<\beta<N$ with $1/p+(N-\beta)/N+1/r=2$. 
Let $u\in L^p(\mathbb{R}^N)$ and $v\in L^r(\mathbb{R}^N)$. 
Then there exists a sharp constant $C(N,p,\beta)$, independent of $u$ and $v$, 
such that
\[
\big|\int_{\mathbb{R}^N}\int_{\mathbb{R}^N}\frac{u(x)v(y)}{|x-y|^{N-\beta}}dxdy\big|
\leq C_{N,p,r,\beta}\| u\|_p \| v\| _r.
\]
If $p=r=\frac{2N}{N+\beta}$, then
\[
 C_{N,p,r,\beta}=C_{N,\beta}
 =\pi^{\frac{N-\beta}{2}}
 \frac{\Gamma(\frac{\beta}{2})}{\Gamma(\frac{N+\beta}{2})}
 \Big\{\frac{\Gamma(\frac{N}{2})}{\Gamma(N)}\Big\}^{-\frac{\beta}{N}}.
\]
\end{lemma}

\begin{lemma}[Gagliardo-Nirenberg inequality of Power type \cite{GN23}] \label{Gn75}
 Let $N\geq1$ and $2\leq p<2^*$, then the following sharp Gagliardo-Nirenberg 
inequality
 \begin{equation}\label{GN}
\| u\|_p\leq C_{N,p}\| u\|_2^{1-\delta_p}\| \nabla u\|_2^{\delta_p}
\end{equation}
holds for any $u\in H^1(\mathbb{R}^N)$, where $\delta_p=\frac{N}{2}-\frac{N}{p}$, 
the sharp constant $C_{N,p}$ is
\[
 C_{N,p}^p=\frac{2p}{2N+(2-N)p}(\frac{2N+(2-N)p}{N(p-2)})^{\frac{N(p-2)}{4}}
 \frac{1}{\|Q_p\|_2^{p-2}}
\]
and $Q_p$ is the unique positive radial solution of the equation
\[
 -\Delta Q+Q=|Q|^{p-2}Q.
\]
\end{lemma}

\begin{lemma}[Gagliardo-Nirenberg inequality of Choquard type\cite{Mza}] \label{123}
Let $N\geq3$, $\alpha\in(0,N)$ and $\frac{N+\alpha}{N}<p<\frac{N+\alpha}{N-2}$, we have
\begin{equation}\label{GN2}
 \int_{\mathbb{R}^N}(I_\alpha\ast|u|^p)|u|^pdx
 \leq C_{N,p}\Big(\int_{\mathbb{R}^N}|\nabla u|
 ^2dx\Big)^{\frac{Np-(N+\alpha)}{2}}
 \Big(\int_{\mathbb{R}^N}|u|^2dx\Big)^{\frac{N+\alpha-p(N-2)}{2}},
\end{equation}
where equality holds for $u=Q_p$, $C_{N,p}=\frac{p}{|Q_p|_2^{2p-2}}$ and $Q_p$ 
is a nontrivial solution of
\begin{equation}
 -\frac{N(p-2)+N-\alpha}{2}\Delta Q_p+\frac{N+\alpha-(N-2)p}{2}Q_p
 =(I_\alpha\ast|Q_p|^p)|Q_p|^{p-2}Q_p.
\end{equation}
\end{lemma}
By \cite[(3.3)]{banqun}, we have the following Lemma.

\begin{lemma}\label{qxz3}
If $\frac{N+\alpha}{N}<r_1,r_2<\frac{N+\alpha}{N-2}$, then
\[
\int_{\mathbb{R}^N}(I_\alpha\ast|u_1|^{r_1})|u_2|^{r_2}dx
\leq\Big(\int_{\mathbb{R}^N}(I_\alpha\ast|u_1|^{r_1})|u_1|^{r_1}dx\Big)^{1/2}
\Big(\int_{\mathbb{R}^N}(I_\alpha\ast|u_2|^{r_2})|u_2|^{r_2}dx\Big)^{1/2},
\]
where $I_\alpha$ is given by \eqref{gamma}.
\end{lemma}


\begin{lemma}[{\cite[Lemma A.2]{em}}] \label{Liuv}
If $p\in(0,\frac{N}{N-2}]$ when $N\geq3$ and $p\in(0,\infty)$ when $N=1,2$. 
Let $u\in L^p(\mathbb{R}^N)$ be a smooth nonnegative function
satisfying $-\Delta u\geq0$ in $\mathbb{R}^N$. Then $u\equiv0$.
\end{lemma}

We now consider the problem
\begin{equation}\label{sig}
\begin{gathered}
-(a+b\int_{\mathbb{R}^N}|\nabla u|^2dx)\Delta u
=\lambda u+\mu(I_\alpha\ast|u|^p)|u|^{p-2}u,\\
\int_{\mathbb{R}^N}|u|^2dx=d>0,
\end{gathered}
\end{equation}
where $a,b>0$, $\frac{N+\alpha}{N}<p<\frac{N+\alpha+2}{N}$ and $N=3,4$. 
The corresponding minimization problem of \eqref{sig} is
\begin{equation}\label{ehs}
\sigma_p^{\mu}(d):=\inf_{u\in S(d)} G_\mu(u),
\end{equation}
where $G_\mu(u)=\frac{a}{2}\| \nabla u\|_2^2+\frac{b}{4}\| \nabla u\|_2^4
-\frac{\mu}{2p} \int_{\mathbb{R}^N}(I_\alpha\ast|u|^p)|u|^pdx$.

\begin{lemma}\label{byg}
\begin{itemize}
\item[(i)]   For any $d>0$, we have $\sigma_p^{\mu}(d)<0$.
\item[(ii)] $\sigma_p^{\mu}(d)$ is continuous with respect to $d\geq0$.
\item[(iii)] For any $d\geq m\geq0$, then 
$\sigma_p^{\mu}(d)\leq\sigma_p^{\mu}(m)+\sigma_p^{\mu}(d-m)$.
\end{itemize}
\end{lemma}

\begin{proof}
Item (i) follows from \cite[Theorem 1.2]{ygg}. 
For item (ii), we assume that $d^n=d+o_n(1)$. By the definition of
$\sigma_p^{\mu}(d^n)$, for any $\varepsilon>0$, there exists $u^n\in
 S_r(d^n)$ such that
\begin{equation}\label{gmb5}
 G_\mu(u^n)\leq\sigma_p^{\mu}(d^n)+\varepsilon.
\end{equation}
Let $w^n:=\frac{u^n}{\| u^n \|_2d^{1/2}}$. It is
easy to check that $w^n\in S_r(d^n)$ and
\begin{equation}\label{gmb6}
 \sigma_p^{\mu}(d)\leq G_\mu(w^n)=G_\mu(u^n)+o_n(1).
\end{equation}
Combining \eqref{gmb5} and \eqref{gmb6}, we have
\[
 \sigma_p^{\mu}(d)\leq\sigma_p^{\mu}(d^n)+\varepsilon+o_n(1).
\]
Reversing the argument, we obtain similarly that
\[
\sigma_p^{\mu}(d^n)\leq\sigma_p^{\mu}(d)+\varepsilon+o_n(1).
\]
Thus, by the arbitrariness of $\varepsilon$, we have
$\sigma_p^{\mu}(d^n)=\sigma_p^{\mu}(d)+o_n(1)$. This proves (ii).

By the density of $C_0^{\infty}(\mathbb{R}^N)$ in $H^1(\mathbb{R}^N)$, 
for any $\varepsilon>0$, there exist $\tilde{\psi}$, $\bar{\psi}
 \in C_0^{\infty}(\mathbb{R}^N)\times C_0^{\infty}(\mathbb{R}^N)$ with 
$\| \tilde{\psi}\|_2^2=m$, $\| \bar{\psi}\|_2^2=d-m$,
$i=1,2$ such that
\begin{gather*}
G_\mu(\tilde{\psi})\leq\sigma_p^{\mu}(m)+\varepsilon,\\
G_\mu(\bar{\psi})\leq\sigma_p^{\mu}(d-m)+\varepsilon.
\end{gather*}
Since $G_\mu$ is invariant by translation, without loss of generality, we may assume 
that $\operatorname{supp}\tilde{\psi}\cap \operatorname{supp}\bar{\psi}=\emptyset$, 
and then
 $\| \tilde{\psi}+\bar{\psi}\|_2^2=\| \tilde{\psi}\|_2^2+\|\bar{\psi}\|_2^2=d$,
\[
 \sigma_p^{\mu}(d)\leq G_\mu(\tilde{\psi}+\bar{\psi})\leq\sigma_p^{\mu}(m)+
 \sigma_p^{\mu}(d-m)+2\varepsilon.
\]
Therefore,
\[
 \sigma_p^{\mu}(d)\leq\sigma_p^{\mu}(m)+\sigma_p^{\mu}(d-m).
\]
This proof is complete.
\end{proof}

\begin{lemma}[{\cite[Lemma 2.5]e{linjie}}] \label{ph2}
Let $\frac{N+\alpha}{N}<r_1,r_2<\frac{N+\alpha}{N-2}$, if
\[
(u_1^n,u_2^n)\rightharpoonup(u_1,u_2) \text{in} H^1(\mathbb{R}^N)\times H^1(\mathbb{R}^N),
\]
then, up to a subsequence,
\begin{equation}\label{Blm}
 \int_{\mathbb{R}^N}(I_\alpha\ast|u_1^n|^{r_1})|u_2^n|^{r_2}dx
 =\int_{\mathbb{R}^N}(I_\alpha\ast|u_1|^{r_1})|u_2|^{r_2}dx+o_n(1).
\end{equation}
\end{lemma}

\begin{lemma}\label{fud}
If {\rm (A1)} holds, then $\inf_{W(k)}F(u_1,u_2)<0$ for all $k>0$.
\end{lemma}

\begin{proof}
Let $u^{\tau}(x)=\tau^{\frac{N}{2}}u(\tau x)$. Then it is easy to check that
\[
(u_1,u_2)\in S_r(d_1)\times S_r(d_2), \quad
 (u_1^{\tau},u_2^{\tau})\in W(k)
\]
when $\tau$ is sufficiently small.
From
\begin{align*}
F(u_1^{\tau},u_2^{\tau})
&= \frac{a}{2}\tau^2\sum_{i=1}^2\| \nabla u_i\|_2^2+\frac{b}{4}\tau^4\sum_{i=1}^2\|
 \nabla u_i\|_2^4-\sum_{i=1}^2\frac{\mu_i}{2p_i}\tau^{N{p_i}-N-\alpha}
  \int_{\mathbb{R}^N}(I_\alpha\ast|u_i|^{p_i})|u_i|^{p_i}dx\\
&\quad -\Theta\tau^{\frac{Nr}{2}-N-\beta}\int_{\mathbb{R}^N}(I_\beta\ast|u_1|^{p_1})
 |u_2|^{p_2}dx,
\end{align*}
$\frac{N+\alpha}{N}<p_1,p_2<\frac{N+\alpha+2}{N}$ and 
$\frac{N+\beta}{N}<r_1,r_2<2_\beta^*$, we have $\Phi_{(u_1,u_2)}(\tau) <0$ for
 $\tau$ small enough.
\end{proof}

\begin{lemma}\label{dx5}
Assume that {\rm (A2)} holds. Then there exist $k_0=k_0(d_1,d_2)>0$, 
$\Theta_*=\Theta_*(d_1,d_2)>0$ such that for any $0<\Theta<\Theta_\star$,
\[
\inf_{W(2k_0)\setminus W(k_0)}F(u_1,u_2)>0.
\]
And there exists $\varepsilon_0>0$ small enough, such that
\[
\sigma(d_1,d_2)<\inf_{\overline{W(k_0)}\setminus A(k_0-\varepsilon_0)}F(u_1,u_2).
\]
\end{lemma}

\begin{proof}
For  $(u_1,u_2)\in S_r(d_1)\times S_r(d_2)$, 
let $k=\| \nabla u_1\|_2^2+\| \nabla u_2\|_2^2$. Then
by Lemma \ref{123} and Lemma \ref{qxz3}, we have
\begin{equation}\label{dmk7}
\begin{aligned}
&F(u_1,u_2) \\
&=\frac{a}{2}k+\frac{b}{4}\sum_{i=1}^2\| \nabla u_i\|_2^4-
\sum_{i=1}^2\frac{\mu_i}{2p_i}\int_{\mathbb{R}^N}
 (I_\alpha\ast|u_i|^{p_i})|u_i|^{p_i}dx-\Theta\int_{\mathbb{R}^N}
 (I_\beta\ast|u_1|^{p_1})|u_2|^{p_2}dx\\
&\geq \frac{b}{8}k^2-\sum_{i=1}^2B_i\| \nabla u_i\|_2^{Np_i-N-\alpha}
 -\Theta B_3k^{\frac{Nr-2(N+\beta)}{4}}\\
&\geq \frac{b}{8}k^2-\sum_{i=1}^2B_ik^{\frac{Np_i-N-\alpha}{2}}
 -\Theta B_3k^{\frac{Nr-2(N+\beta)}{4}}:=g(k),\\
\end{aligned}
\end{equation}
where $B_i=\frac{\mu_i}{2p_i}C_{N,p_i}d_i^{\frac{N+\alpha-p_i(N-2)}{2}}$ ($i=1,2$), 
and $B_3=C_{N,r_1}C_{N,r_2}d_1^{\frac{N+\beta-r_1(N-2)}{4}}
d_2^{\frac{N+\beta-r_2(N-2)}{4}}$.
By (A2), we have
\[
0<\frac{Np_i-N-\alpha}{2}<1 \quad \text{and}\quad
  \frac{Nr-2(N+\beta)}{4}>2, \quad i=1,2.
\]
We fix $k=k_0$ sufficiently large such that
\[
\sum_{i=1}^2B_ik_0^{\frac{Np_i-N-\alpha-4}{2}}\leq\frac{b}{24},
\]
and fix $\Theta=\Theta_*$ sufficiently small such that
\[
\Theta_*B_3(2k_0)^{\frac{Nr-2(N+\beta)-8}{4}}\leq\frac{b}{24}.
\]
Therefore, for each $0<\Theta\leq\Theta_*$ and
$(u_1,u_2)\in W(2k_0)\setminus W(k_0)$,
we obtain
\begin{align*}
F(u_1,u_2)
&\geq \frac{b}{8}k^2-\sum_{i=1}^2B_ik^{\frac{Np_i-N-\alpha}{2}}
 -\Theta B_3k^{\frac{Nr-2(N+\beta)}{4}}\\
&= k^2\Big(\frac{b}{8}-\sum_{i=1}^2B_ik^{\frac{Np_i-N-\alpha-4}{2}}
 -\Theta B_3k^{\frac{Nr-2(N+\beta)-8}{4}}\Big)\\
&\geq bk_0^2\Big(\frac{1}{8}-\frac{1}{24}-\frac{1}{24}\Big)
=\frac{b}{24}k_0^2.
\end{align*}

Next, by continuity of $g(k)$ and $g(k_0)>0$, there exists $\varepsilon_0>0$ 
sufficiently small such that $g(k)\geq0$ when $k\in[k_0-\varepsilon,k_0]$. Hence,
\[
F(u_1,u_2)\geq g(k)\geq0>\sigma(d_1,d_2)
\]
for any $(u_1,u_2)\in\overline{W(k_0)}\setminus W(k_0-\varepsilon_0)$.
This proof is complete.
\end{proof}

\begin{lemma}\label{zjngz}
Let $\frac{N+\alpha}{N}<p_1,p_2<2_\alpha^*$ and $\frac{N+\beta}{N}<r_1,r_2<2_\beta^*$. 
If $(u_1,u_2)\neq(0,0)$ is a solution of
 \eqref{yy} for some $(\lambda_1,\lambda_2)\in\mathbb{R}^2$, then $\lambda_1<0$ 
 or $\lambda_2<0$. Furthermore, if $u_1>0$ and $u_2\geq0$, then
$\lambda_1<0$; if $u_1\geq0$ and $u_2>0$, then $\lambda_2<0$.
\end{lemma}

\begin{proof}
Testing \eqref{yy} by $(u_1,u_2)$ and integrating in $\mathbb{R}^N$, one has
\begin{align*}
&\lambda_1\| u_1\|_2^2+\lambda_2\| u_2\|_2^2 \\
&=a\sum_{i=1}^2\| \nabla u_i\|_2^2+
b\sum_{i=1}^2\| \nabla u_i\|_2^4-\sum_{i=1}^2\mu_i\int_{\mathbb{R}^N}(I_\alpha\ast|u_i|^{p_i})|u_i|^{p_i}dx-
\Theta r\int_{\mathbb{R}^N}(I_\beta\ast|u_1|^{r_1})|u_2|^{r_2}dx.
\end{align*}
Combining $P(u_1,u_2)=0$, $2<p_1,p_2<2^*$ and $\frac{N+\beta}{N}<r_1,r_2<2_\beta^*$,
we have
 \begin{align*}
 \lambda_1\| u_1\|_2^2+\lambda_2\| u_2\|_2^2
&= \sum_{i=1}^2\mu_i(Np_i-N-\alpha-1)
 \int_{\mathbb{R}^N}(I_\alpha\ast|u_i|^{p_i})|u_i|^{p_i}dx\\
&\quad +\Theta\left(\frac{Nr-2(N+\beta)}{4}-r\right)\int_{\mathbb{R}^N}(I_\beta\ast|u_1|^{p_1})|u_2|^{p_2}dx<0.
 \end{align*}
Thus, at least one of $\lambda_1$ and $\lambda_2$ is negative.

Then, we argue by contradiction, suppose $\lambda_1\geq0$. 
In view $u_1>0$, we have
\[
-\Big(a+b\int_{\mathbb{R}^N}|\nabla u_1|^2dx\Big)\Delta u_1
=\lambda_1u_1+\mu_1(I_\alpha\ast|u_1|^{p_1})|u_1|^{p_1-2}u_1
 +\Theta r_1(I_\beta\ast|u_2|^{p_2})|u_1|^{p_1-2}u_1\geq0.
\]
By Lemma \ref{Liuv}, we obtain $u_1=0$, which  contradicts $u_1\neq0$.
Therefore, $\lambda_1<0$. The other case can be proved
in the same way.
\end{proof}

\begin{lemma}\label{zyy4}
Assume that $2<p_1,p_2<\frac{2N}{N-2}$ and 
$\frac{N+\beta}{N}<r_1,r_2<\frac{N+\beta}{N-2}$. For any bounded Palais-Smale sequence
$\{(u_1^n,u_2^n)\}$ for $F$ on $S_r(d_1)\times S_r(d_2)$, then there exist 
$(u_1,u_2)\in H_r^1(\mathbb{R}^N)\times H_r^1(\mathbb{R}^N)$ and a sequence
$\{(\lambda_1^n,\lambda_2^n)\}\subset\mathbb{R}^2$, such that up to a subsequence
\begin{itemize}
\item[(a)] $(u_1^n,u_2^n)\rightharpoonup (u_1,u_2)$ in 
 $H_r^1(\mathbb{R}^N)\times H_r^1(\mathbb{R}^N)$, 
 $(u_1^n,u_2^n)\to  (u_1,u_2)$ in $L^p
(\mathbb{R}^N)\times L^p(\mathbb{R}^N)$ for $p\in(2,\frac{2N}{N-2})$.

\item[(b)] $(\lambda_1^n,\lambda_2^n) \to  (\lambda_1,\lambda_2)$ in
 $(\mathbb{R}^2)$.
 
\item[(c)] $F'(u_1^n,u_2^n)-\lambda_1^n(u_1^n,0)-\lambda_2^n(0,u_2^n)\to 0$
 in $H_r^{-1}(\mathbb{R}^N)\times H_r^{-1}(\mathbb{R}^N)$.
 
\item[(d)] $(u_1^n,u_2^n)$ is a solution to the system \eqref{yy} for some 
 $\lambda_1,\lambda_2\leq0$ if $(u_1,u_2)$ satisfies the additional property
$P(u_1^n,u_2^n)\to 0$, where $(\lambda_1,\lambda_2)$ is given by (b).
\end{itemize}
\end{lemma}

\begin{proof}
Obviously, item (a) is true 
Since $\{(u_1^n,u_2^n)\}\subset H_r^1(\mathbb{R}^N)\times H_r^1(\mathbb{R}^N)$ 
is bounded, by \cite{Arm2},
we have $(F|_{S_r(d_1)\times S_r(d_2)})'(u_1^n,u_2^n)\to 0$ in 
$H_r^{-1}(\mathbb{R}^N)\times H_r^{-1}(\mathbb{R}^N)$ is
equivalent to
\[
F'(u_1^n,u_2^n)-\frac{1}{\| u_1^n\|_2^2}\langle F'(u_1^n,u_2^n),
(u_1^n,0)\rangle(u_1^n,0)-
\frac{1}{\| u_2^n\|_2^2}\langle F'(u_1^n,u_2^n),(0,u_2^n)\rangle(0,u_2^n)\to 0
\]
in $H_r^{-1}(\mathbb{R}^N)\times H_r^{-1}(\mathbb{R}^N)$. Thus, we have
\[
 F'(u_1^n,u_2^n)-\lambda_1^n(u_1^n,0)-\lambda_2^n(0,u_2^n)\to 0
 \quad \text{in } H_r^{-1}(\mathbb{R}^N)\times H_r^{-1}(\mathbb{R}^N)
\]
with
\begin{equation}\label{yjj1}
\begin{aligned}
\lambda_1^n&=\frac{1}{\| u_1^n\|_2^2}\Big(a\| \nabla u_1^n\|_2^2+b\| \nabla u_1^n\|_2^4-
\mu_1\int_{\mathbb{R}^N}(I_\alpha\ast|u_1^n|^{p_1})|u_1^n|^{p_1}dx \\
&\quad -\Theta r_1\int_{\mathbb{R}^N}(I_\beta\ast|u_1^n|^{p_1})|u_2^n|^{p_2}dx\Big)
-o_n(1),
\end{aligned}
\end{equation}
and
\begin{equation}\label{yjj2}
\begin{aligned}
 \lambda_2^n
&=\frac{1}{\| u_2^n\|_2^2}\Big(a\| \nabla u_2^n\|_2^2+b\| \nabla u_2^n\|_2^4-
 \mu_2\int_{\mathbb{R}^N}(I_\alpha\ast|u_2^n|^{p_2})|u_2^n|^{p_2}dx \\
&\quad -\Theta r_2\int_{\mathbb{R}^N}(I_\beta\ast|u_1^n|^{p_1})|u_2^n|^{p_2}dx\Big)-o_n(1).
\end{aligned}
 \end{equation}
This proves (c).
By the boundedness of $u_1^n,u_2^n$ in $H_r^1(\mathbb{R}^N)$, $L^p(\mathbb{R}^N)$
for $p\in(2,2^*)$, Lemma \ref{123} and Lemma \ref{qxz3},
we obtain $\{\lambda_1^n\}$, $\{\lambda_2^n\}$ are bounded. This proves (b).
Combining (b) and (c), it is now standard to deduce (d).
\end{proof}

\begin{lemma}\label{yuk9}
Under the conditions of Lemma \ref{zyy4},  $u_1^n\to  u_1$ in $H_r^1(\mathbb{R}^N)$
if $\lambda_1<0$. Similarly,
$u_2^n\to  u_2$ in $H_r^1(\mathbb{R}^N)$ if $\lambda_2<0$.
\end{lemma}

\begin{proof}
In view of Lemmas \ref{ph2} and  \ref{zyy4}, we obtain
\begin{gather}\label{vbp1}
\int_{\mathbb{R}^N}(I_\alpha\ast|u_i^n|^{p_i})|u_i^n|^{p_i}dx
\to \int_{\mathbb{R}^N}(I_\alpha\ast|u_i|^{p_i})|u_i|^{p_i}dx 
(\frac{N+\alpha}{N}<p_i<2_\alpha^*,\quad  i=1,2), \\
\label{vbp2}
\int_{\mathbb{R}^N}(I_\beta\ast|u_1^n|^{r_1})|u_2^n|^{r_2}dx
\to \int_{\mathbb{R}^N}(I_\beta\ast|u_1|^{r_1})|u_2|^{r_2}dx (\frac{N+\beta}{N}
<r_1,r_2<2_\beta^*), \\
\langle F'(u_1^n,u_2^n)-\lambda_1^n(u_1^n,0),(u_1^n,0)\rangle\to 0
= \langle F'(u_1,u_2)-\lambda_1^n(u_1,0),(u_1,0)\rangle.
\end{gather}
Therefore,
\[
a\| \nabla u_1^n\|_2^2+b\| \nabla u_1^n\|_2^4-\lambda_1^n\| u_1^n\|_2^2\to
a\| \nabla u_1\|_2^2+b\| \nabla u_1\|_2^4-\lambda_1\| u_1\|_2^2.
\]
Since
\begin{equation}\label{vbp3}
 \| u_1\|_2^2\leq \varliminf_{n\to +\infty} \| u_1^n\|_2^2, \quad
 \| \nabla u_1\|_2^2\leq \varliminf_{n\to +\infty} \| \nabla u_1^n\|_2^2,
\end{equation}
it follows that
\[
 \| \nabla u_1^n\|_2^2\to \| \nabla u_1\|_2^2, \quad
 \| u_1^n\|_2^2\to \| u_1\|_2^2.
\]
This proof is complete.
\end{proof}

\section{Proof of Theorem \ref{thm1.1}} 

\begin{lemma}\label{xj5}
Assume that {\rm (A1)} holds. Then $F$ is bounded form below and coercive on 
$S_{r}(d_1)\times S_{r}(d_2)$. Furthermore, there exists a bounded
Palais-Smale sequence $\{(u_1^n,u_2^n)\}\subset S_{r}(d_1)\times S_{r}(d_2)$, 
which satisfies $(u_1^n)^-\to 0$ and
$(u_2^n)^-\to 0$ in $H_r^1(\mathbb{R}^N)\times H_r^1(\mathbb{R}^N)$.
\end{lemma}

\begin{proof}
Since $\frac{N+\alpha}{N}<p_1,p_2<\frac{N+\alpha+2}{N}$, 
$2\cdot\frac{N+\beta}{N}<r_1+r_2<2\cdot\frac{N+\beta+4}{N}$, we have $0<Np_i-N-\alpha<2$ 
and $\frac{Nr-2(N+\beta)}{4}<2$.
From \eqref{dmk7}, we know that $F$ is bounded form below and coercive on 
$S_{r}(d_1)\times S_{r}(d_2)$.

Next, let $\{(u_1^n,u_2^n)\}\subset S_{r}(d_1)\times S_{r}(d_2)$ be a minimizing 
sequence for $F$. By the coerciveness of $F$, the sequence is
bounded. Since the functional $F$ is even, we may assume that $w_1^n\geq0$ and 
$w_2^n\geq0$. It is easy to check that $F$ is a $C^1-$manifold in
$H_r^1(\mathbb{R}^N)\times H_r^1(\mathbb{R}^N)$, then by Ekeland's variational 
principle, there exists a minimizing sequence
$\{(u_1^n,u_2^n)\}\subset S_{r}(d_1)\times S_{r}(d_2)$, which is the Palais-Smale 
sequence for $F$ restricted to $S_{r}(d_1)\times S_{r}(d_2)$ and satisfies
$\| (u_1^n,u_2^n)-(w_1^n,w_2^n)\|_{H_r^1(\mathbb{R}^N)\times H_r^1(\mathbb{R}^N)}\to 0$ 
as $n\to +\infty$.
From $w_1^n\geq0$ and $w_2^n\geq0$, we have $(u_1^n)^-\to 0$ and $(u_2^n)^-\to 0$ in 
$H_r^1(\mathbb{R}^N)\times H_r^1(\mathbb{R}^N)$.
\end{proof}

\begin{proof}[Proof of Theorem \ref{thm1.1}] 
By Lemma \ref{xj5}, there exists a bounded Palais-Smale sequence 
$\{(u_1^n,u_2^n)\}\subset S_{r}(d_1)\times S_{r}(d_2)$
and $(u_1,u_2)\in H_r^1(\mathbb{R}^N)\times H_r^1(\mathbb{R}^N)$ with $u_1\geq0$ 
and $u_2\geq0$, such that $\lim_{n\to +\infty}
(u_1^n,u_2^n)=\sigma(d_1,d_2)$, $(u_1^n,u_2^n)\rightharpoonup(u_1,u_2)$ in 
$H_r^1(\mathbb{R}^N)\times H_r^1(\mathbb{R}^N)$. Hence, it suffices to
prove that $(u_1^n,u_2^n)\to (u_1,u_2)$ in 
$H_r^1(\mathbb{R}^N)\times H_r^1(\mathbb{R}^N)$. Indeed, if this holds, then we have
$\{(u_1^n,u_2^n)\}\subset S_{r}(d_1)\times S_{r}(d_2)$ and 
$F(u_1,u_2)=\sigma(d_1,d_2)$. Furthermore, by the strong maximum principle,
we have $u_1>0$ and $u_2>0$. From $\Theta>0$, obviously, we have 
$\sigma(d_1,d_2)\leq\sigma_{p_1}^{\mu_1}(d_1)+\sigma_{p_2}^{\mu_2}(d_2)$. 
By Lemma \ref{fud}, we know
$\sigma(d_1,d_2)<0$. We divide four cases to have that $u_1>0$ and $u_2>0$.
\smallskip 

\noindent\textbf{Case I:} $u_1=0$ and $u_2=0$.
Obviously, we can obtain that $\lim_{n\to +\infty}F(u_1^n,u_2^n)=0$, which is 
 contradicts $\sigma(d_1,d_2)<0$.
\smallskip

\noindent\textbf{Case II:} $u_1=0$ and $u_2\neq0$.
Combining \eqref{vbp1} and \eqref{vbp3}, we have
\begin{equation}\label{kgf}
\varliminf_{n\to +\infty}F(u_1^n,u_2^n)
\geq\frac{a}{2}\| \nabla u_2\|_2^2+\frac{b}{4}\|
 \nabla u_2\|_2^4-\frac{\mu_2}{2p_2}
 \int_{\mathbb{R}^N}(I_\alpha\ast|u_2|^{p_2})|u_2|^{p_2}dx
\geq \sigma_{p_2}^{\mu_2}(d_3),
\end{equation}
where $d_3:=\| u_2\|_2^2\leq d_2$. By (i) and (iii) of Lemma \ref{byg}, 
we have $\sigma_{p_2}^{\mu_2}(d_3)\geq\sigma_{p_2}^{\mu_2}(d_2)$.
Thus, by \eqref{kgf}, we have $\sigma(d_1,d_2)\geq \sigma_{p_2}^{\mu_2}(d_3)$. 
On the other hand, using Lemma \ref{byg} again, we have $\sigma_{p_1}^{\mu_1}(d_1)<0$. 
In view of $\sigma(d_1,d_2)\leq\sigma_{p_1}^{\mu_1}(d_1)+\sigma_{p_2}^{\mu_2}(d_2)$, 
we have $\sigma(d_1,d_2)<\sigma_{p_2}^{\mu_2}(d_2)$.
This is a contradiction.
\smallskip

\noindent\textbf{Case III:} $u_1\neq0$ and $u_2=0$.
By the same proof as in Case II, we can  obtain a contradiction.
\smallskip

\noindent\textbf{Case IV:} $u_1\neq0$ and $u_2\neq0$.
From Lemma \ref{zjngz}, we have $\lambda_1<0$ and $\lambda_2<0$. Thus, by 
Lemma \ref{yuk9}, we obtain $u_1^n\to  u_1$ and
$u_2^n\to  u_2$ in $H_r^1(\mathbb{R}^N)$. This proof is complete.
\end{proof}

\section{Proof of Theorem \ref{thm1.2}}

Let $k=\| \nabla u_1\|_2^2+\| \nabla u_2\|_2^2$, then for each 
$(u_1,u_2)\in S_r(d_1)\times S_r(d_2)$, by the same
proof as in Lemma \ref{dx5}, we obtain
\begin{align*}
&F(u_1,u_2) \\
&= \frac{a}{2}\sum_{i=1}^2\| \nabla u_i\|_2^2
 +\frac{b}{4}\sum_{i=1}^2\| \nabla u_i\|_2^4
 - \sum_{i=1}^2\frac{\mu_i}{2p_i}\int_{\mathbb{R}^N}(I_\alpha\ast|u_i|^{p_i})
 |u_i|^{p_i}dx-\Theta\int_{\mathbb{R}^N}(I_\beta\ast|u_1|^{p_1})|u_2|^{p_2}dx\\
&\geq \frac{b}{8}k^2-B_1k^{\frac{Np_1-N-\alpha}{2}}-B_2k^{\frac{Np_2-N-\alpha}{2}}
 -\Theta B_3k^{\frac{Nr-2(N+\beta)}{4}}:=g(k),
\end{align*}
where $B_i=\frac{\mu_i}{2p_i}C_{N,p_i}d_i^{\frac{N+\alpha-p_i(N-2)}{2}}$ ($i=1,2$), 
and $B_3=C_{N,r_1}C_{N,r_2}d_1^{\frac{N+\beta-r_1(N-2)}{4}}
d_2^{\frac{N+\beta-r_2(N-2)}{4}}$.

By Lemmas \ref{fud} and \ref{dx5}, we consider the  minimization problem
\[
\sigma(d_1,d_2):=\inf_{W(k_0)} F(u_1,u_2)<0 \text{for any} 0<\Theta\leq\Theta_*.
\]

\begin{lemma}\label{ddh3}
Assume that {\rm (A1)} holds and $0<\alpha-2\leq\beta<\alpha<N$. Then
there exists $\Theta_1>0$, such that if $0<\Theta<\Theta_1$, then the function $g(k)$
 has a unique local minimum point at the negative
level and a unique global maximum point at the positive level.
 Moreover, there exists $0<k_0<k_1$, such that $g(k_0)=g(k_1)=0$ and $g(k)>0$ 
 if and only if $k\in(k_0,k_1)$.
\end{lemma}

\begin{proof}
Without loss of generality, we assume that $p_1\geq p_2$.
\smallskip

\noindent\textbf{Case I:} $p_1=p_2=p$. For $k>0$, we have
 \begin{align*}
 g(k)&= \frac{b}{8}k^2-(B_1+B_2)k^{\frac{Np-N-\alpha}{2}}
  -\Theta B_3k^{\frac{Nr-2(N+\beta)}{4}}\\
 &= k^{\frac{Np-N-\alpha}{2}}\Big(\frac{b}{8}k^{\frac{4+N+\alpha-Np}{2}}
  -(B_1+B_2)-\Theta B_3k^{\frac{Nr-2Np+2(\alpha-\beta)}{4}}\Big).
 \end{align*}
Define $h(k)=\frac{b}{8}k^{\frac{4+N+\alpha-Np}{2}}
-\Theta B_3k^{\frac{Nr-2Np+2(\alpha-\beta)}{4}}$, then 
$g(k)>0\Leftrightarrow h(k)>B_1+B_2$.
Setting
\[
\tilde{k}=\left(\frac{b(4+N+\alpha-Np)}{4\Theta B_3(Nr-2Np
+2(\alpha-\beta))}\right)^{\frac{4}{Nr-2(N+\beta+4)}},
\]
it is easy to check that $h(k)$ is increasing in $(0,\tilde{k})$, 
decreasing in $(\tilde{k},+\infty)$ and
 \[
 h(\tilde{k})=M\Theta^{-\frac{2(4+N+\alpha-Np)}{Nr-2(N+\beta+4)}},
 \]
where $M$ is a positive constant.
Hence, we can take $\Theta$ small enough, such that $h(\tilde{k})>B_1+B_2$.
Thus, there exists $\Theta_1>0$, such that
$0<\Theta<\Theta_1$, which implies that $g(k)>0$ on $k\in(k_0,k_1)$.
Since $g(k)\to 0^-$ as $k\to 0^+$, $g(k)$ has a local minimum point on
$(0,k_0)$. Therefore, $g(k)$ has at least two critical points.

For $k>0$, we define
\[
 \varphi(k)=\frac{b}{4}k^{\frac{4+N+\alpha-Np}{2}}
 -\frac{Nr-2(N+\beta)}{4}\Theta B_3k^{\frac{Nr-2Np+2(\alpha-\beta)}{4}}.
\]
By direct computation,
$g'(k)=k^{\frac{Np-N-\alpha-2}{2}}(\varphi(k)-(B_1+B_2)(Np-N-\alpha))$ and
$g'(k)=0\Leftrightarrow \varphi(k)=(B_1+B_2)(Np-N-\alpha)$.
Since $\varphi(k)$ has a unique global maximum point,
 $\varphi(k)=(B_1+B_2)(Np-N-\alpha)$ has at most two solutions, i.e., $g(k)$ has at
most two critical points. Thus, $g(k)$ has a unique local minimum
point at the negative level and a unique global maximum point at the positive level.
\smallskip

\noindent\textbf{Case II:} $p_1>p_2$. In this case, we have
 \begin{align*}
 g(k)
 &= \frac{b}{8}k^2-B_1k^{\frac{Np_1-N-\alpha}{2}}-B_2k^{\frac{Np_2-N-\alpha}{2}}-\Theta B_3k^{\frac{Nr-2(N+\beta)}{4}}\\
 &= k^{\frac{Np_2-N-\alpha}{2}}\Big(\frac{b}{8}k^{\frac{4+N+\alpha-Np_2}{2}}
  -B_1k^{\frac{Np_1-Np_2}{2}}-B_2-
 \Theta B_3k^{\frac{Nr-2Np_2+2(\alpha-\beta)}{4}}\Big)\\
&=:k^{\frac{Np_2-N-\alpha}{2}}(Q(k)-B_2).
 \end{align*}
By calculation and analysis, $Q(k)$ has a unique global maximum point $\bar{k}$ with
\[
 \bar{k}>\Big(\frac{b(4+N+\alpha-Np_1)(4+N+\alpha-Np_2)}{2B_3\Theta(Nr-2Np_1
 +2(\alpha-\beta))(Nr-2Np_2+2(\alpha-\beta))}\Big)^{\frac{4}{Nr-2(N+\beta+4)}}:=k_*,
\]
and $Q(\bar{k})>Q(k_*)>{\bar{k}}^{\frac{N(p_2-p_1)}{2}}$ for $\Theta$ small enough.
Therefore, we can take $\Theta$ small enough,
such that $Q(\bar{k})>B_2$. The rest of the proof is similar to the case $p_1=p_2$.
\end{proof}

\begin{lemma}\label{zlx}
Assume that {\rm (A1)} holds and $0<\alpha-2\leq\beta<\alpha<N$. Then there exists 
$\Theta_2>0$, such that if $0<\Theta<\Theta_2$, then $\mathcal{P}_0(d_1,d_2)=\emptyset$, 
and $\mathcal{P}$ is a submanifold
of $H_r^1(\mathbb{R}^N)\times H_r^1(\mathbb{R}^N)$.
\end{lemma}

\begin{proof}
We argue by contradiction. Suppose that there exists
 $(u_1,u_2)\in\mathcal{P}_0(d_1,d_2)$. Then
\begin{equation}\label{paza24}
\begin{aligned}
&a\sum_{i=1}^2\| \nabla u_i\|_2^2+b\sum_{i=1}^2\| \nabla u_i\|_2^4
-\sum_{i=1}^2\frac{\mu_i}{2p_i}(Np_i-N-\alpha)
\int_{\mathbb{R}^N}(I_\alpha\ast|u_i|^{p_i})|u_i|^{p_i}dx\\
&-\Theta\Big(\frac{N}{2}-N-\beta\Big)\int_{\mathbb{R}^N}(I_\beta\ast|u_1|^{r_1})|u_2|^{r_2}dx=0
\end{aligned}
\end{equation}
and
\begin{equation}\label{paza25}
\begin{aligned}
&2a\sum_{i=1}^2\| u_i\|_2^2+4b\sum_{i=1}^2\| u_i\|_2^4
-\sum_{i=1}^2\frac{\mu_i}{2p_i}(Np_i-N-\alpha)^2
\int_{\mathbb{R}^N}(I_\alpha\ast|u_i|^{p_i})|u_i|^{p_i}dx\\
&-\Theta\Big(\frac{Nr}{2}-N-\beta\Big)^2\int_{\mathbb{R}^N}
 (I_\beta\ast|u_1|^{r_1})|u_2|^{r_2}dx=0.
\end{aligned}
\end{equation}
By \eqref{paza24} and \eqref{paza25}, we have
\begin{equation}\label{tnh3}
 \begin{aligned}
&\Big(\frac{Nr}{2}-N-\beta-2\Big)a\sum_{i=1}^2\| u_i\|_2^2
+(\frac{Nr}{2}-N-\beta-4)b\sum_{i=1}^2\| u_i\|_2^4\\
&=\sum_{i=1}^2\frac{\mu_i}{2p_i}\Big(\frac{Nr}{2}-Np_i+\alpha-\beta\Big)
\int_{\mathbb{R}^N}(I_\alpha\ast|u_i|^{p_i})|u_i|^{p_i}dx.
\end{aligned}
\end{equation}
On the one hand, by \eqref{tnh3}, we have
\begin{align*}
\frac{b}{2}(\| \nabla u_1\|_2^2+\| \nabla u_2\|_2^2)^2&\leq b(\| \nabla u_1\|_2^4
+\| \nabla u_2\|_2^4)+\frac{\frac{Nr}{2}-N-\beta-2}{\frac{Nr}{2}-N-\beta-4}
 a\sum_{i=1}^2\| \nabla u_i\|_2^2\\
&= \frac{1}{\frac{Nr}{2}-N-\beta-4}\sum_{i=1}^2\frac{\mu_i}{2p_i}
 \Big(\frac{Nr}{2}-Np_i+\alpha-\beta\Big)
 \int_{\mathbb{R}^N}(I_\alpha\ast|u_i|^{p_i})|u_i|^{p_i}dx\\
&\leq \sum_{i=1}^2\mathbb{K}_i(\| \nabla u_1\|_2^2
 +\| \nabla u_2\|_2^2)^{\frac{Np_i-N-\alpha}{2}},
\end{align*}
where 
\[
\mathbb{K}_i:=\frac{1}{\frac{Nr}{2}-N-\beta-4}\frac{\mu_i}{2p_i}
\Big(\frac{Nr}{2}-Np_i+\alpha-\beta\Big)C_{N,p_i,\alpha}^{p_i}
d_i^{\frac{N+\alpha-p_i(N-2)}{2}}.
\]
Thus, there exists $\mathbb{K}_3>0$, such that
\begin{equation}\label{hand1}
\| \nabla u_1\|_2^2+\| \nabla u_2\|_2^2\leq\mathbb{K}_3.
\end{equation}
On the other hand, combining \eqref{paza24} and \eqref{paza25}, we have
\begin{align*}
&2b\sum_{i=1}^2\| u_i\|_2^4+\sum_{i=1}^2\frac{\mu_i}{2p_i}
(Np_i-N-\alpha)(N+\alpha+2-Np_i)\int_{\mathbb{R}^N}(I_\alpha\ast
 |u_i|^{p_i})|u_i|^{p_i}dx\\
&= \Theta\Big(\frac{Nr}{2}-N-\beta\Big)\Big(\frac{Nr}{2}-N-\beta-2\Big)
 \int_{\mathbb{R}^N}(I_\beta\ast|u_1|^{r_1})|u_2|^{r_2}dx.
\end{align*}
Thus,
\begin{equation}\label{yhg3}
\begin{aligned}
&\sum_{i=1}^2\frac{\mu_i}{2p_i}(Np_i-N-\alpha)(N+\alpha+2-Np_i)\int_{\mathbb{R}^N}(I_\alpha\ast|u_i|^{p_i})|u_i|^{p_i}dx\\
&\leq 2b\sum_{i=1}^2\| u_i\|_2^4
+\sum_{i=1}^2\frac{\mu_i}{2p_i}(Np_i-N-\alpha)(N+\alpha+2-Np_i)\int_{\mathbb{R}^N}(I_\alpha\ast|u_i|^{p_i})|u_i|^{p_i}dx\\
&= \Theta\Big(\frac{Nr}{2}-N-\beta\Big)
\Big(\frac{Nr}{2}-N-\beta-2\Big)\int_{\mathbb{R}^N}(I_\beta\ast|u_1|^{r_1})|u_2|^{r_2}dx.
\end{aligned}
\end{equation}
If $\int_{\mathbb{R}^N}(I_\beta\ast|u_1|^{r_1})|u_2|^{r_2}dx=0$, by \eqref{yhg3}, 
we have $u_1=u_2=0$, which is a contradiction.
By \eqref{yhg3}, we obtain
\begin{align*}
&\frac{b}{2}(\| u_1\|_2^2+\| u_2\|_2^2)^2\\
&\leq b(\| u_1\|_2^4+\| u_2
\|_2^4)+a(\| u_1\|_2^2+\| u_2\|_2^2)\\
&= \sum_{i=1}^2\frac{\mu_i}{2p_i}(Np_i-N-\alpha)
\int_{\mathbb{R}^N}(I_\alpha\ast|u_i|^{p_i})|u_i|^{p_i}dx+\Theta
\Big(\frac{Nr}{2}-N-\beta\Big)\int_{\mathbb{R}^N}(I_\beta\ast|u_1|^{r_1})|u_2|^{r_2}dx\\
&\leq \Theta\Big(\frac{Nr}{2}-N-\beta\Big)
\Big(\frac{\frac{Nr}{2}-N-\beta-2}{N+\alpha+2-N\cdot\min\{p_1,p_2\}}+1\Big)
\int_{\mathbb{R}^N}(I_\beta\ast|u_1|^{r_1})|u_2|^{r_2}dx\\
&\leq \Theta\mathbb{K}_4(\| \nabla u_1\|_2^2+\| \nabla u_2\|_2^2)^{\frac{Nr-2(N+\beta)}{4}}.
\end{align*}
Since $\frac{rN-2(N+\beta)}{4}>2$, then we have
\begin{equation}\label{hand2}
 \| \nabla u_1\|_2^2+\| \nabla u_2\|_2^2
 \geq\Big(\frac{b}{2\mathbb{K}_4\Theta}\Big)^{\frac{4}{Nr-2(N+\beta+4)}}.
\end{equation}
Combining \eqref{hand1} and \eqref{hand2}, we obtain
\[
 \Big(\frac{b}{2\mathbb{K}_4\Theta}\Big)^{\frac{4}{Nr-2(N+\beta+4)}}
 \leq\| \nabla u_1\|_2^2+\| \nabla u_2\|_2^2
 \leq \mathbb{K}_3;
\]
it does not hold for $\Theta$ small enough.

Next, we show that $\mathcal{P}(d_1,d_2)$ is a submanifold of 
$H_r^1(\mathbb{R}^N)\times H_r^1(\mathbb{R}^N)$. Note that
\[
\mathcal{P}(d_1,d_2):=\{(u_1,u_2)\in H_r^1(\mathbb{R}^N)
\times H_r^1(\mathbb{R}^N):P(u_1,u_2)=0, \mathbb{S}_1(u_1)=0,\mathbb{S}_2(u_2)=0\},
\]
where $\mathbb{S}_1(u_1)=\| u_1\|_2^2-d_1$, $\mathbb{S}_2(u_2)=\| u_2\|_2^2-d_2$.
It suffices to show that
$d(P,\mathbb{S}_1,\mathbb{S}_2):H_r^1(\mathbb{R}^N)\times H_r^1(\mathbb{R}^N)
\to \mathbb{R}^3$ is surjective. Otherwise, by the
independence of $d\mathbb{S}_1(u_1)$ and $d\mathbb{S}_2(u_2)$, $dP(u_1,u_2)$ must
be a linear combination of $d\mathbb{S}_1(u_1)$ and $d\mathbb{S}_2(u_2)$,
i.e., there exists $\kappa_1,\kappa_2\in\mathbb{R}$, such that $(u_1,u_2)$
is a weak solution of
\begin{align*}
&-(a+2b\int_{\mathbb{R}^N}|\nabla u_1|^2dx)\Delta u_1 \\
&= \kappa_1 u_1+\mu_1\frac{Np_1-N-\alpha}{2}(I_\alpha\ast|u_1|^{p_1})|u_1|^{p_1-2}u_1
 +\Theta r_1\frac{Nr-2(N+\beta)}{4}(I_\beta\ast|u_2|^{p_2})|u_1|^{p_1-2}u_1,
\\
&-(a+2b\int_{\mathbb{R}^N}|\nabla u_2|^2dx)\Delta u_2 \\
&= \kappa_2 u_2+\mu_2\frac{Np_2-N-\alpha}{2}(I_\alpha\ast|u_2|^{p_2})|u_2|^{p_2-2}u_2
 +\Theta r_2\frac{Nr-2(N+\beta)}{4}(I_\beta\ast|u_1|^{p_1})|u_2|^{p_2-2}u_2. \\
\end{align*}
Testing the above system with $(u_1,u_2)$ and combining with Poho\v{z}aev identity,
we obtain $(\Phi_{(u_1,u_2)})''(0)=0$.
Then $(u_1,u_2)\in\mathcal{P}_0(d_1,d_2)=0$, which is contradicted to
$\mathcal{P}_0(d_1,d_2)=\emptyset$. The proof is complete.
\end{proof}

Let $\Theta_*=\min\{\Theta_1,\Theta_2\}$,
$\varPi_1:=\{(u_1,u_2)\in H_r^1(\mathbb{R}^N)\times 
H_r^1(\mathbb{R}^N):\int_{\mathbb{R}^N}(I_\beta
\ast|u_1|^{r_1})|u_2|^{r_2}dx>0\}$ and 
$\varPi_2:=\{(u_1,u_2)\in H_r^1(\mathbb{R}^N)\times H_r^1(\mathbb{R}^N):
\int_{\mathbb{R}^N}(I_\beta \ast|u_1|^{r_1})|u_2|^{r_2}dx=0\}$. 
Next, we analyze the geometric structure of $\Phi_{(u_1,u_2)}(\tau)$.

\begin{lemma}\label{threep}
Assume that {\rm (A1)} holds and let $0<\Theta<\Theta_*$. 
Then for each $(u_1,u_2)\in S_r(d_1)\times S_r(d_2)\cap\varPi_1$, 
$\Phi_{(u_1,u_2)}(\tau)$
has exactly two critical points $\tau_{(u_1,u_2)}<s_{(u_1,u_2)}$ and two zeros
 $e_{(u_1,u_2)}<f_{(u_1,u_2)}$ with $\tau_{(u_1,u_2)}<e_{(u_1,u_2)}<s_{(u_1,u_2)}<
f_{(u_1,u_2)}$. Moreover,
\begin{itemize}
\item[(i)] $\tau\star(u_1,u_2)\in\mathcal{P}_+(d_1,d_2)$ if and only if 
$\tau=\tau_{(u_1,u_2)}$, $\tau\star(u_1,u_2)\in\mathcal{P}_-(d_1,d_2)$ if and only if 
$\tau=s_{(u_1,u_2)}$.

\item[(ii)] If $\tau<e_{(u_1,u_2)}$, then 
$\| \nabla(\tau\star u_1)\|_2^2+\| \nabla(\tau\star u_2)\|_2^2\leq k_0$, and
$F(\tau_{(u_1,u_2)}\star(u_1,u_2))=\min\{F(\tau\star(u_1,u_2)):
\tau\in\mathbb{R},\| \nabla(\tau\star u_1)\|_2^2+\|
\nabla(\tau\star u_2)\|_2^2<k_0\}<0$. Moreover, 
$F(s_{(u_1,u_2)}\star(u_1,u_2))=\max_{\tau\in\mathbb{R}}F(\tau\star(u_1,u_2))$.

\item[(iii)] The maps $(u_1,u_2)\mapsto \tau_{(u_1,u_2)}$ and 
$(u_1,u_2)\mapsto s_{(u_1,u_2)}$ are of class $C^1$.
\end{itemize}
\end{lemma}

\begin{proof}
For  $(u_1,u_2)\in S_r(d_1,d_2)$, we obtain
\begin{align*}
(\Phi_{(u_1,u_2)})'(\tau)
&= ae^{2\tau}\sum_{i=1}^2\| \nabla u_i\|_2^2+be^{4\tau}\sum_{i=1}^2\| \nabla u_i\|_2^4
 -\sum_{i=1}^2\frac{\mu_i}{2p_i}e^{(Np_i-N-\alpha)\tau}\int_{\mathbb{R}^N}(I_\alpha\ast|u_i|^{p_i})|u_i|^{p_i}dx\\
&\quad -\Theta(\frac{Nr}{2}-N-\beta) e^{(\frac{Nr}{2}-N-\beta)\tau}
 \int_{\mathbb{R}^N}(I_\beta\ast|u_1|^{r_1})|u_2|^{r_2}dx\\
&= a\sum_{i=1}^2\| \nabla (\tau\star u_i)\|_2^2+b\sum_{i=1}^2
 \| \nabla (\tau\star u_i)\|_2^4 \\
&\quad -\sum_{i=1}^2 \frac{\mu_i}{2p_i}(Np_i-N-\alpha)
\int_{\mathbb{R}^N}(I_\alpha\ast|\tau\star u_i|^{p_i})|\tau\star u_i|^{p_i}dx\\
&\quad -\Theta(\frac{Nr}{2}-N-\beta)\int_{\mathbb{R}^N}(I_\beta\ast|\tau\star u_1|^{r_1})
|\tau\star u_2|^{r_2}dx\\
&= P(\tau\star u_1,\tau\star u_2).
\end{align*}
Therefore, $\tau\star(u_1,u_2)\in\mathcal{P}(d_1,d_2)$ if and only if 
$(\Phi_{(u_1,u_2)})'(\tau)=0$. Obviously,
\[
\Phi_{(u_1,u_2)}(\tau)=F(\tau\star u_1,\tau\star u_2)\geq g(e^{2\tau}k).
\]
By Lemma \ref{ddh3}, we have
\[
 \Phi_{(u_1,u_2)}(\tau)>0, \quad \text{for all }
 \tau\in\Big(\frac{1}{2}\ln\frac{k_0}{k},\frac{1}{2}\ln\frac{k_1}{k}\Big).
\]
In view of $\Phi_{(u_1,u_2)}(-\infty)=0^-$ and $\Phi_{(u_1,u_2)}(+\infty)=-\infty$,
it follows that $\Phi_{(u_1,u_2)}(\tau)$ has at least two critical points
$\tau_{u_1,u_2}<s_{(u_1,u_2)}$, where $s_{(u_1,u_2)}$ is global maximum point at
positive level, $\tau_{(u_1,u_2)}$ is a local minimum point on
$(-\infty,\frac{1}{2}\ln\frac{k_0}{k})$ at negative level.
Similar to Lemma \ref{ddh3}, $\Phi_{(u_1,u_2)}(\tau)$ has at most two critical points.
Therefore, $\Phi_{(u_1,u_2)}(\tau)$ has exactly the two critical points
$\tau_{(u_1,u_2)}$ and $s_{(u_1,u_2)}$.
From $\tau\star(u_1,u_2)\in\mathcal{P}(d_1,d_2)$ if and only
$\Phi_{(u_1,u_2)}'(\tau)=0$, then $\tau\star(u_1,u_2)\in\mathcal{P}(d_1,d_2)$
if and only $\tau=\tau_{(u_1,u_2)}$ or $\tau=s_{(u_1,u_2)}$.
Since $\tau_{(u_1,u_2)}$ is a local minimum point, we have that
\[
(\Phi_{\tau_{(u_1,u_2)}\star(u_1,u_2)})''(0)=(\Phi_{(u_1,u_2)}''(\tau_{(u_1,u_2)})\geq0.
\]
Since $\mathcal{P}_0(d_1,d_2)=\emptyset$, it follows that
 $(\Phi_{\tau_{(u_1,u_2)}\star(u_1,u_2)})''(0)>0$.
Then $\tau_{(u_1,u_2)}\star(u_1,u_2)\in\mathcal{P}_+(d_1,d_2)$
and $s_{(u_1,u_2)}\star(u_1,u_2)\in\mathcal{P}_-(d_1,d_2)$.
By the monotonicity, $\Phi_{(u_1,u_2)}(\tau)$ has exactly two zeros
$e_{(u_1,u_2)}<f_{(u_1,u_2)}$ with
$\tau_{(u_1,u_2)}<e_{(u_1,u_2)}<s_{(u_1,u_2)}<f_{(u_1,u_2)}$.
Finally, let $\varPsi(\tau,u_1,u_2):=(\Phi_{(u_1,u_2)})'(\tau)$. Then
\begin{gather*}
\varPsi(\tau_{(u_1,u_2)},u_1,u_2)=\varPsi(s_{(u_1,u_2)},u_1,u_2)=0\\
\partial_{\tau}\varPsi(\tau_{(u_1,u_2)},u_1,u_2)=(\Phi_{(u_1,u_2)})''(\tau_{(u_1,u_2)})>0,\\
\partial_{s}\varPsi(s_{(u_1,u_2)},u_1,u_2)=(\Phi_{(u_1,u_2)})''(s_{(u_1,u_2)})<0.\\
\end{gather*}
Applying the implicit function theorem, the maps $(u_1,u_2)\mapsto \tau_{(u_1,u_2)}$
and $(u_1,u_2)\mapsto s_{(u_1,u_2)}$ are of class $C^1$.
\end{proof}

\begin{lemma}\label{jiap6}
Assume that {\rm (A1)} holds and let $0<\Theta<\Theta_*$. Then for each
$(u_1,u_2)\in S_r(d_1)\times S_r(d_2)\cap\varPi_2$, $\Phi_{(u_1,u_2)}(\tau)$
 has a unique critical point $\tau_{(u_1,u_2)}$ and a zero $e_{(u_1,u_2)}$ with 
$\tau_{(u_1,u_2)}<e_{(u_1,u_2)}$. Moreover,
\begin{itemize}
\item[(i)] $\mathcal{P}(d_1,d_2)=\mathcal{P}_+(d_1,d_2)$ and 
 $\tau\star(u_1,u_2)\in\mathcal{P}_+(d_1,d_2)$ if and only if $\tau=\tau_{(u_1,u_2)}$.
 
\item[(ii)] $F(\tau_{(u_1,u_2)}\star(u_1,u_2))=\min_{\tau\in\mathbb{R}}
 F(\tau\star(u_1,u_2))$.
 
\item[(iii)] $\| \nabla(\tau\star u_1)\|_2^2+\| \nabla(\tau\star u_2)\|_2^2<k_0$ for all 
$\tau<e_{(u_1,u_2)}$.
\end{itemize}
 \end{lemma}

\begin{proof}
Without loss of generality, we assume that $p_1\geq p_2$. It is easy to see that 
$\Phi_{(u_1,u_2)}(\tau)\to 0^-$ as $\tau\to -\infty$,
and $\Phi_{(u_1,u_2)}(\tau)\to +\infty$ as $\tau\to +\infty$ for any 
$(u_1,u_2)\in S_r(d_1,d_2)\cap\varPi_2$. Hence, $\Phi_{(u_1,u_2)}(\tau)$ has
a global minimum point $\tau_{(u_1,u_2)}$ at negative level. Note that 
$(\Phi_{(u_1,u_2)})'(\tau)=0$ if and only if
\begin{align*}
&ae^{(N+\alpha+2-Np_2)\tau}\sum_{i=1}^2\| \nabla u_i\|_2^2
 +be^{(N+\alpha+4-Np_2)\tau}\sum_{i=1}^2\| \nabla u_i\|_2^4\\
&-\frac{\mu_1}{2p_1}(Np_1-N-\alpha)e^{N(p_1-p_2)\tau}
 \int_{\mathbb{R}^N}(I_\alpha\ast|u_1|^{p_1})|u_1|^{p_1}dx \\
&=\frac{\mu_2}{2p_2}(Np_2-N-\alpha)\int_{\mathbb{R}^N}(I_\alpha\ast|u_2|^{p_2})
 |u_2|^{p_2}dx.
\end{align*}
By direct calculation, it is not difficult to check that equation has exactly one 
solution. Hence, $\tau\star(u_1,u_2)\in\mathcal{P}(d_1,d_2)$
if and only if $\tau=\tau_{(u_1,u_2)}$. Since $\tau_{(u_1,u_2)}$ is global minimum 
point of $\Phi_{(u_1,u_2)}(\tau)$, then $(\Phi_{(u_1,u_2)})''(\tau_{(u_1,u_2)})\geq0$.
Since $\mathcal{P}_0(d_1,d_2)=\emptyset$, we obtain that 
$(\Phi_{\tau_{(u_1,u_2)}\star(u_1,u_2)})''(0)>0$, and thus 
$\tau_{(u_1,u_2)}\star(u_1,u_2)\in\mathcal{P}_+(d_1,d_2)$.
Furthermore, using the monotonicity and the behavior at infinity, 
$\Phi_{(u_1,u_2)}(\tau)$ has a unique zero $e_{(u_1,u_2)}$ with 
$\tau_{(u_1,u_2)}<e_{(u_1,u_2)}$.
By $\Phi_{(u_1,u_2)}(\tau)\geq g(e^{2\tau}(\| \nabla u_1\|_2^2+\| \nabla u_2\|_2^2))$, 
then $\Phi_{(u_1,u_2)}(\tau)\geq g(k_0)=0$ at 
$\tau=\frac{1}{2}\ln\frac{k_0}{\| \nabla u_1\|_2^2+\| \nabla u_2\|_2^2}$.
Therefore, $\| \nabla(\tau\star u_1)\|_2^2+\| \nabla(\tau\star u_2)\|_2^2<k_0$ 
for all $\tau<e_{(u_1,u_2)}$.
\end{proof}

\begin{lemma}\label{xiajy2}
Assume that {\rm (A1)} holds. If $0<\Theta<\Theta_*$, then
\[
\sigma(d_1,d_2)=\inf_{\mathcal{P}(d_1,d_2)}F(u_1,u_2)=\inf_{\mathcal{P}_+(d_1,d_2)}F(u_1,u_2),
\]
and there exists $\varepsilon_0>0$ small enough, such that
\[
\sigma(d_1,d_2)<\inf_{\overline{W(k_0)} \setminus W(k_0-\varepsilon_0)}F(u_1,u_2).
\]
\end{lemma}

\begin{proof}
We first prove that $\sigma(d_1,d_2)=\inf_{\mathcal{P}_+(d_1,d_2)}F(u_1,u_2)$.
For any $(u_1,u_2)\in\mathcal{P}_+(d_1,d_2)$, $\tau_{(u_1,u_2)}=0$. 
By Lemma \ref{threep},
$0<\frac{1}{2}\ln\frac{k_0}{\| \nabla u_1\|_2^2+\| \nabla u_2\|_2^2}$,
so $k_0>\| \nabla u_1\|_2^2+
\| \nabla u_2\|_2^2$. Hence, $\mathcal{P}_+(d_1,d_2)\subset W(k_0)$ and
\[
\sigma(d_1,d_2)\leq\inf_{\mathcal{P}_+(d_1,d_2)}F(u_1,u_2).
\]
On the other hand, for any $(u_1,u_2)\in W(k_0)$, there exists a unique
$\tau_{(u_1,u_2)}\in\mathbb{R}$, such that
\[
\tau_{(u_1,u_2)}\star(u_1,u_2)\in\mathcal{P}_+\subset W(k_0).
\]
Using (ii) in Lemma \ref{threep} and Lemma \ref{jiap6}, we have
\[
F(\tau_{(u_1,u_2)}\star(u_1,u_2))=\min\{F(\tau\star(u_1,u_2)):\tau\in\mathbb{R},\| \nabla(\tau\star u_1)\|_2^2+\|
\nabla(\tau\star u_2)\|_2^2<k_0\}\leq F(u_1,u_2).
\]
Hence, $\sigma(d_1,d_2)\geq\inf_{\mathcal{P}_+(d_1,d_2)}F(u_1,u_2)$. Therefore,
$\sigma(d_1,d_2)=\inf_{\mathcal{P}_+(d_1,d_2)}F(u_1,u_2)$.
By (ii) in Lemma \ref{threep}, we obtain
\[
 \inf_{\mathcal{P}(d_1,d_2)}F(u_1,u_2)=\inf_{\mathcal{P}_+(d_1,d_2)}F(u_1,u_2).
\]
This proof is complete.
By $g(k_0)=0$ and continuity of $g$, there is $\varepsilon_0>0$ such that
$g(k)\geq\frac{\sigma(d_1,d_2)}{2}$ if $k\in[k_0-\varepsilon,k_0]$. Therefore,
\[
F(u_1,u_2)\geq g(k)\geq\frac{\sigma(d_1,d_2)}{2}>\sigma(d_1,d_2) \quad
\text{for all } (u_1,u_2)\in\overline{W(k_0)}\setminus W(k_0-\varepsilon_0).
\]
\end{proof}

\begin{lemma}\label{bijia}
Assume that {\rm(A1)} holds. Let $0<\Theta<\Theta_*$, then
$\sigma(d_1,d_2)<\min\{\sigma_0(d_1,0),\sigma_0(d_2,0)\}$.
\end{lemma}

\begin{proof}
For each $(u_1,u_2)\in W(k_0)$, we have
\begin{align*}
F(u_1,u_2)
&= G_{\mu_1}(u_1)+G_{\mu_2}(u_2)-\Theta\int_{\mathbb{R}^N}(I_\beta\ast|u_1|^{r_1})|u_2|^{r_2}dx\\
&\leq G_{\mu_1}(u_1)+G_{\mu_2}(u_2).
\end{align*}
Hence, $\sigma(d_1,d_2)\leq\inf_{W(k_0)}(G_{\mu_1}(u_1)+G_{\mu_2}(u_2))$.
 For any $u_1\in S_r(d_1)$ with $\| \nabla u_1\|_2^2=k_0$, then
\begin{align*}
G_{\mu_1}(u_1)
&= \frac{a}{2}\| \nabla u_1\|_2^2+\frac{b}{4}\| \nabla u_1\|_2^4-
\frac{\mu_1}{p_1}\int_{\mathbb{R}^N}(I_\alpha\ast|u_1|^{2p_1})|u_1|^{p_1}dx\\
&\geq \frac{b}{4}k_0^2-B_i\| \nabla u_1\|_2^{\delta_{p_i}p_i}\geq g(k_0)=0,
\end{align*}
where $B_i$ is defined by Lemma \ref{dx5}. Therefore,
\[
\inf_{S_r(d_1)}G_{\mu_1}=\inf_{E(d_1,k_0)}G_{\mu_1}(u_1)<0,
\]
where $E(d,k):=\{u\in S_r(d):\| \nabla u\|_2^2<k\}$. Since the map
$k\mapsto\frac{b}{4}k_0^2-B_i\| \nabla u_1\|_2^{\delta_{p_i}p_i}$
is continuity, by similar proof in Lemma \ref{xiajy2}, there is
$\varepsilon_0>0$ such that
\[
\sigma(d_1,\mu_1)<\inf_{\overline{E(d_1,k_0)}
\setminus E(d_1,k_0-\varepsilon_0)}G_{\mu_1}(u_1),
\]
and $\inf_{E(d_2,\varepsilon_0)}G_{\mu_2}(u_2)<0$.
We define $\Lambda:=\{(u_1,u_2):u_1\in E(d_1,k_0-\varepsilon_0),u_2
\in E(d_2,\varepsilon_0)\}$; then
$\Lambda\subset W(k_0)$. Therefore,
\begin{align*}
\inf_{W(k_0)}(G_{\mu_1}(u_1)+G_{\mu_2}(u_2))
&\leq \inf_{\Lambda}(G_{\mu_1}(u_1)+G_{\mu_2}(u_2))\\
&= \inf_{E(d_1,k_0-\varepsilon_0)}G_{\mu_1}(u_1)
 +\inf_{E(d_2,\varepsilon_0)}G_{\mu_2}(u_2)\\
&<\inf_{E(d_1,k_0-\varepsilon_0)}G_{\mu_1}(u_1)\\
&= \sigma_0(d_1,0).
\end{align*}
Thus, $\sigma(d_1,d_2)<\sigma_0(d_1,0)$. It can be shown that
$\sigma(d_1,d_2)<\sigma_0(d_2,0)$ in the same way as above.
\end{proof}

\begin{lemma}\label{zuih9}
Let $\{(u_1^n,u_2^n)\}\subset S_r(d_1)\times S_r(d_2)$ be a Palais-Smale 
sequence for $F|_{S_r(d_1)\times S_r(d_2)}$ at level $\sigma_{\beta}(d_1,d_2)$, 
and $P(u_1^n,u_2^n) \to 0$ as $n\to +\infty$. Then, $(u_1^n,u_2^n)\to (u_1,u_2)$ 
in $H_r^1(\mathbb{R}^N)\times H_r^1(\mathbb{R}^N)$.
\end{lemma}

\begin{proof}
By $F(|u_1^n|,|u_2^n|)=F(u_1^n,u_2^n)$, then we can assume that $u_1^n,u_2^n\geq0$. 
Since $\{(u_1^n,u_2^n)\}\subset S_r(d_1)\times S_r(d_2)$ be a
Palais-Smale sequence for $F|_{S_r(d_1)\times S_r(d_2)}$ at level $\sigma(d_1,d_2)$, 
we have
\begin{align*}
 \sigma(d_1,d_2)+o_n(1)
&= \frac{a}{2}\sum_{i=1}^2\| \nabla u_i^n\|_2^2+\frac{b}{4}\sum_{i=1}^2\| \nabla u_i^n\|_2^4-
 \sum_{i=1}^2\frac{\mu_i}{2p_i}\int_{\mathbb{R}^N}(I_\alpha\ast|u_i^n|^{p_i})|u_i^n|^{p_i}dx\\
&\quad -\Theta\int_{\mathbb{R}^N}(I_\beta\ast|u_1^n|^{p_1})|u_2^n|^{p_2}dx\\
&= a\left(\frac{1}{2}-\frac{2}{Nr-2(N+\beta)}\right)\sum_{i=1}^2\| \nabla u_i^n\|_2^2
 +b\left(\frac{1}{4}-\frac{2}{Nr-2(N+\beta)}\right)\sum_{i=1}^2\| \nabla u_i^n\|_2^4\\
&\quad -\sum_{i=1}^2\frac{\mu_i}{2p_i}\left(1-\frac{2(Np_i-N-\alpha)}{Nr-2(N+\beta)}\right)\int_{\mathbb{R}^N}(I_\alpha\ast|u_i^n|^{p_i})|u_i^n|^{p_i}dx\\
&\geq a\left(\frac{1}{2}-\frac{2}{Nr-2(N+\beta)}\right)k+\frac{b}{2}\left(\frac{1}{2}-\frac{2}{Nr-2(N+\beta)}\right)k^2\\
&\quad-\sum_{i=1}^2\frac{\mu_i}{2p_i}M_i\left(1-\frac{2(Np_i-N-\alpha)}{Nr-2(N+\beta)}\right)\|\nabla u_i^n\|_2^{Np_i-N-\alpha},\\
\end{align*}
where $k=\| \nabla u_1^n\|_2^2+\| \nabla u_2^n\|_2$,
$M_i:=C_{N,p_i}^{p_i}d_i^{\frac{N+\alpha-(N-2)p_i}{2}}$ ($i=1,2$),
$8<Nr-2(N+\beta)<4\cdot2_\beta^*$ and $0<Np_i-N-\alpha<2$. 
We obtain the sequence
$\{(u_1^n,u_2^n)\}$ is bounded in $H_r^1(\mathbb{R}^N)\times H_r^1(\mathbb{R}^N)$. 
Hence, we have
\begin{gather*}
(u_1^n,u_2^n)\rightharpoonup (u_1,u_2) \quad  \text{in } 
H_r^1(\mathbb{R}^N)  \times H_r^1(\mathbb{R}^N)\\
(u_1^n,u_2^n)\to  (u_1,u_2) \quad \text{in } 
L^p(\mathbb{R}^N)\times L^p(\mathbb{R}^N) \text{for} p\in(2,2^*).
\end{gather*}
By (c) in Lemma \ref{zyy4}, there exist two sequences of real numbers 
$\{\lambda_1^n\}$, $\{\lambda_2^n\}$, such that
\begin{align*}
&a\sum_{i=1}^2\int_{\mathbb{R}^N}\nabla u_i^n\nabla\phi_idx+b\sum_{i=1}^2\int_{\mathbb{R}^N}|\nabla u_i^n|^2dx\int_{\mathbb{R}^N}\nabla u_i^n\nabla\phi_idx\\
&-\sum_{i=1}^2\mu_i\int_{\mathbb{R}^N}(I_\alpha\ast|u_i^n|^{p_i})|u_i^n|^{p_i-2}u_i^n\phi_idx-\Theta r_1\int_{\mathbb{R}^N}(I_\beta\ast|u_2^n|^{r_2})|u_1^n|^{r_1-2}u_1^n\phi_1dx\\
&-\Theta r_2\int_{\mathbb{R}^N}(I_\beta\ast|u_1^n|^{r_1})|u_2^n|^{r_2-2}u_2^n\phi_1dx-\sum_{i=1}^2\int_{\mathbb{R}^N}\lambda_i^nu_i^n\phi_idx\\
&= o_n(1)\| (\phi_,\phi_2)\|_{H_r^1(\mathbb{R}^N)\times H_r^1(\mathbb{R}^N)},
\end{align*}
where $o_n(1)\to 0$ as $n\to +\infty$. By Lemma \ref{zyy4}, $\lambda_i^n\to \lambda_i$
 $(i=1,2)$, and $(u_1,u_2)$ is a solution to \eqref{yy}.
Since
\begin{equation}\label{ddk3}
F(u_1,u_2)\leq\varliminf_{n\to +\infty}F(u_1^n,u_2^n)=\sigma(d_1,d_2)<0,
\end{equation}
we have $(u_1,u_2)\neq(0,0)$. From Lemma \ref{zjngz}, we obtain
that $\lambda_1<0$ or $\lambda_2<0$. Without loss of generality, we assume that 
$\lambda_1<0$. Using Lemma \ref{yuk9}, $u_1^n\to  u_1$ in $H_r^1(\mathbb{R}^N)$.
Suppose by contradiction that $\lambda_2\geq0$. Then
\[
 -(a+b\int_{\mathbb{R}^N}|\nabla u_2|^2dx)\Delta u_2=\lambda_2 u_2+\mu_2(I_\alpha\ast|u_2|^{p_2})|u_2|^{p_2-2}u_2+
 \Theta r_2(I_\beta\ast|u_1|^{r_1})|u_2|^{r_2-2}u_2\geq0.
\]
By Lemma \ref{Liuv}, we have $u_2=0$. And $u_1$ satisfies
\begin{gather*}
-(a+b\int_{\mathbb{R}^N}|\nabla u_1|^2dx)\Delta u
=\lambda_1 u_1+\mu_1(I_\alpha\ast|u_1|^{p_1})|u_1|^{p_1-2}u_1,\\
\int_{\mathbb{R}^N}|u|^2dx=d_1>0.
 \end{gather*}
Hence,
\begin{align*}
\sigma_{\beta}(d_1,d_2)
&= \lim_{n\to +\infty} F(u_1^n,u_2^n)\\
&= \lim_{n\to +\infty} \frac{a}{2}\sum_{i=1}^2\| \nabla u_i\|_2^2+\frac{b}{4}\sum_{i=1}^2\| \nabla u_i\|_2^4-
\sum_{i=1}^2\frac{\mu_i}{2p_i}\int_{\mathbb{R}^N}(I_\alpha\ast|u_i|^{p_i})|u_i|^{p_i}dx\\
&\quad -\Theta\int_{\mathbb{R}^N}(I_\beta\ast|u_1|^{r_1})|u_2|^{r_2}dx\\
&= \lim_{n\to +\infty} \frac{a}{2}\sum_{i=1}^2\| \nabla u_i\|_2^2+\frac{b}{4}\sum_{i=1}^2\| \nabla u_i\|_2^4
-\frac{\mu_1}{2p_1}\int_{\mathbb{R}^N}(I_\alpha\ast|u_1|^{p_1})|u_1|^{p_1}dx\\
&\geq \frac{a}{2}\| \nabla u_1\|_2^2+\frac{b}{4}\| \nabla u_1\|_2^4-\frac{\mu_1}{2p_1}
\int_{\mathbb{R}^N}(I_\alpha\ast|u_1|^{p_1})|u_1|^{p_1}dx\\
&= G_{\mu_1}(d_1)\\
&\geq  \sigma(d_1,0),
\end{align*}
which  contradicts  Lemma \ref{bijia}. Hence, $\lambda_2<0$, and then,
$u_2^n\to  u_2$ in $H_r^1(\mathbb{R}^N)$.
\end{proof}

\begin{proof}[Proof of Theorem \ref{thm1.2}]
Let $\{(u_1^n,u_2^n)\}\subset W(k_0)$ be a minimizing sequence for
$\sigma(d_1,d_2)$, i.e.\\ $F(u_1^n,u_2^n)\to \sigma(d_1,d_2)$. 
By Lemma \ref{threep}, $\tau_{(u_1^n,u_2^n)}\star(u_1^n,u_2^n)
\in\mathcal{P}_+(d_1,d_2)$ for every $n$, 
$\| \nabla(\tau_{(u_1^n,u_2^n)}\star u_1^n)
\|_2^2+\| \nabla(\tau_{(u_1^n,u_2^n)}\star u_2^n)\|_2^2<k_0$ and
\[
F(\tau_{(u_1^n,u_2^n)}\star u_1^n,\tau_{(u_1^n,u_2^n)}\star u_2^n)\leq F(u_1^n,u_1^n).
\]
Let $(\varpi_1^n,\varpi_2^n):=(\tau_{(u_1^n,u_2^n)}\star u_1^n,\tau_{(u_1^n,u_2^n)}
\star u_2^n)$, then $\{\varpi_1^n,\varpi_2^n\}\subset W(k_0)$ is a minimizing
sequence for $\sigma_{\beta}(d_1,d_2)$ and
$(\varpi_1^n,\varpi_2^n)\in\mathcal{P}_+(d_1,d_2)$. By Lemma \ref{xiajy2},
$\{(\varpi_1^n,\varpi_2^n)\}\subset W(k_0-\varepsilon_0)$.
Therefore, by Ekeland's variational principle, there is a radial symmetric
Palais-Smale sequence $(\upsilon_1^n,\upsilon_2^n)$ for $F|_{S_r(d_1)\times S_r(d_1)}$
satisfying
$\| (\varpi_1^n,\varpi_2^n)-(\upsilon_1^n,\upsilon_2^n)
\|_{H_r^1(\mathbb{R}^N)\times H_r^1(\mathbb{R}^N)}\to 0$ as
$n\to +\infty$. Thus, $\{(\upsilon_1^n,\upsilon_2^n)\}\subset W(k_0)$ and
$P(\upsilon_1^n,\upsilon_2^n)\to 0$ as $n\to +\infty$.
Now Lemma \ref{zuih9} implies that
there exists $\upsilon_1,\upsilon_2>0$ such that
$(\upsilon_1^n,\upsilon_2^n)\to  (\upsilon_1,\upsilon_2)$ in
$H_r^1(\mathbb{R}^N)\times H_r^1(\mathbb{R}^N)$, and then
$(\upsilon_1,\upsilon_2)$ is a local minimizer for $F|_{W(k_0)}$.
Therefore, $(\upsilon_1,\upsilon_2)$ is a positive radial solution to \eqref{yy}
for some $\lambda_1,\lambda_2<0$.
\end{proof}

\section{Proof of Theorem \ref{thm1.3}}

In this section, we prove the existence of the second normalized solution. 
By \cite{Wyu}, the  Choquard equation
\begin{equation}\label{Wmc4}
-\Delta u+u=\int_{\mathbb{R}^N}\frac{|u(y)|^2}{|x-y|^{N-2}}dyu  \quad    \text{in } 
  \mathbb{R}^3
\end{equation}
has a unique positive solution, which is often a strong help to obtain the second 
solution of \eqref{yy}.

By Lemma \ref{fud} and Lemma \ref{dx5}, we introduce a minimax structure of
 mountain pass type. There exists
$k^*\in(0,k_0)$, such that for any $0<\Theta\leq\Theta_*$
\[
\delta(s_1,d_2):=\inf_{h\in\Gamma}\max_{t\in[0,1]}F(h(t))
>\max\{F(h(0)),F(h(1))\},
\]
where $\Gamma:=\{h\in C([0,1],S_r(d_1)\times S_r(d_2)):h(0)
\in W(\bar{k}),h(1)\notin\overline{W(k_0)},F(h(1))<0\}$.

\begin{lemma}\label{str4}
 Assume that {\rm (A2)} holds. Then for each $0<\Theta<\Theta_*$, there exists 
a Palais-Smale sequence $\{(u_1^n,u_2^n)\}$ for $F|_{S_r(d_1)\times S_r(d_2)}$
at the level $\delta(d_1,d_2)$, which satisfies $(u_1^n)^-\to 0$, $(u_2^n)^-\to 0$ 
and $P(u_1^n,u_2^n)\to 0$.
\end{lemma}

\begin{proof}
We recall the stretched functional first introduced in \cite{97Na},
 \[
\tilde{F}: \mathbb{R}\times \Big(H_r^1(\mathbb{R}^N)\times H_r^1(\mathbb{R}^N)\Big)
\to \mathbb{R}, \quad (s,(u_1,u_2))\mapsto F(s\star u_1,s\star u_2).
 \]
We define
\begin{align*}
\tilde{\Gamma}:= \big\{&\tilde{h}\in C([0,1],S_r(d_1)\times S_r(d_2)):\tilde{h}(0)
\in(0,h(0)),\tilde{h}(1)=(0,h(1)),\\
&h(0)\in W(k^*),h(1)\notin\overline{W(k_0)},F(h(1))<0\big\}
\end{align*}
and
\[
\tilde{\delta}(d_1,d_2)
:=\inf_{\tilde{h}\in\tilde{\Gamma}}\max_{t\in[0,1]}\tilde{F}(\tilde{h}(t)).
\]
It is not difficult to check that $\tilde{\delta}(d_1,d_2)=\delta(d_1,d_2)$.
Indeed, by the definitions of $\tilde{\delta}(d_1,d_2)$ and $\delta(d_1,d_2)$,
this identity follows immediately from the fact that the maps
\[
\phi: \Gamma\to \tilde{\Gamma}, \quad h\mapsto \phi(h):(0,h),
\]
and
\[
 \varphi : \tilde{\Gamma}\to \Gamma, \quad
 \tilde{h}=(\varpi,h)\mapsto\varphi(\tilde{h}):=\varpi\star h \quad
  \text{with } (\varpi\star h)(t)=\varpi(t)\star h(t)
\]
satisfying
\[
\tilde{F}(\phi(h))=F(h),\quad  F(\varphi(\tilde{h}))=\tilde{F}(\tilde{h}).
\]
Then, we obtain a sequence $\{(v_1^n,v_2^n)\}\subset\Gamma$ such that
\[
\max_{t\in[0,1]}\tilde{F}(0,(v_1^n(t),v_2^n(t)))\to \tilde{\delta}(d_1,d_2).
\]
In view of $F(u_1,u_2)=F(|u_1|,|u_2|)$ for $(u_1,u_2)\in S_r(d_1)\times S_r(d_2)$,
then we can assume that $v_1^n(t)\geq0$ and $v_2^n(t)\geq0$ for $t\in[0,1]$.
By \cite[Theorem 4.1]{sben}, there exists a Palais-Smale sequence
$\{(s_n,(u_1^n,u_2^n))\}$ for $\tilde{F}|_{\mathbb{R}\times(S_r(d_1)\times S_r(d_2))}$
at the level $\delta(d_1,d_2)$, such that $s_n\to 0$ and
$\| (u_1^n,u_2^n)-(v_1^n,v_2^n)\|\to 0$. It follows that
$(u_1^n)^-\to 0$, $(u_1^n)^-\to 0$.
Since $\tilde{F}(s,(u_1,u_2))=\tilde{F}(0,s\star(u_1,u_2))$, we have
\[
(\partial_s\tilde{F})(s,(u_1,u_2))=(\partial_s\tilde{F})(0,s\star(u_1,u_2))
\]
and $(\partial_u\tilde{F})(s,u)[\psi]=(\partial_u\tilde{F})(0,s\star u)[s\star\psi]$
for $u=(u_1,u_2)$, $\psi=(\psi_1,\psi_2)$. Therefore,
$\{(0,s_n\star(u_1^n,u_2^n))\}$ is also a Palais-Smale sequence for
$\tilde{F}|_{\mathbb{R}\times(S_r(d_1)\times S_r(d_2))}$ at the level
$\delta(d_1,d_2)$.
Then we may assume that $s_n=0$, which implies that
$\{(u_1^n,u_2^n)\}\subset S_r(d_1)\times S_r(d_2)$ is a Palais-Smale sequence for
$F|_{S_r(d_1)\times S_r(d2)}$ at the level $\delta(d_1,d_2)$ and
$(\partial_s\tilde{F})(0,(u_1^n,u_2^n))\to 0$, that $P(u_1^n,u_2^n)\to 0$ holds.
\end{proof}

\begin{lemma}\label{lmk}
Assume that {\rm (A1)} holds and $0<\Theta\leq\Theta_*$, then there exists a 
positive and radial solution $(u_1,u_2)$ to the system \eqref{yy} for some
$(\lambda_1,\lambda_2)$, and $F(u_1,u_2)=\delta(d_1,\delta_2)$.
\end{lemma}

\begin{proof}
By proof similar to that  of Lemma \ref{zuih9}, the Palais-Smale sequence 
$\{(u_1^n,u_2^n)\}$ for $F$ restricted to
$S_r(d_1)\times S_r(d_2)$ at the level $\delta(d_1,d_2)$ is bounded in 
$H_r^1(\mathbb{R}^N)\times H_r^1(\mathbb{R}^N)$.
Now we can assume that
\begin{gather*}
(u_1^n,u_2^n)\rightharpoonup(u_1,u_2) \quad \text{in }
 H_r^1(\mathbb{R}^N)\times H_r^1(\mathbb{R}^N), \\
(u_1^n,u_2^n)\to (u_1,u_2) \quad \text{in } L^p(\mathbb{R}^N)\times 
L^p(\mathbb{R}^N) \text{ for } p\in(2,2^*).
\end{gather*}
From Lemma \ref{zyy4}, there exists a sequence 
$\{(\lambda_1^n,\lambda_2^n)\}\subset\mathbb{R}^2$, such that 
$(\lambda_1^n,\lambda_2^n)\to (\lambda_1,\lambda_2)$ in $\mathbb{R}^2$, 
$(u_1,u_2)$ is a solution to the system \eqref{yy} and $P(u_1,u_2)=0$. 
By $(u_1^n)^-\to 0$, $(u_2^n)^-\to 0$, then $u_1,u_2\geq0$.

Next, we prove that $F(u_1,u_2)=\delta(d_1,d_2)$. By $P(u_1^n,u_2^n)\to 0$, i.e.,
\begin{equation}\label{fmb}
\begin{aligned}
a\sum_{i=1}^2\| \nabla u_i^n\|_2^2+b\sum_{i=1}^2\| \nabla u_i^n\|_2^4 
&= \sum_{i=1}^2\frac{\mu_i}{2p_i}(Np_i-N-\alpha)\int_{\mathbb{R}^N}
 (I_\alpha\ast|u_i^n|^{p_i})|u_i^n|^{p_i}dx \\
&\quad +\Theta\frac{Nr-2(N+\beta)}{2}\int_{\mathbb{R}^N}
(I_\beta\ast|u_1^n|^{r_1})|u_2^n|^{p_2} dx.
\end{aligned}
\end{equation}
By Lemmas \ref{ph2} and \ref{zyy4}, we have
\begin{align*}
&\sum_{i=1}^2\frac{\mu_i}{2p_i}(Np_i-N-\alpha)\int_{\mathbb{R}^N}
 (I_\alpha\ast|u_i^n|^{p_i})|u_i^n|^{p_i}dx+\Theta\frac{Nr-2(N+\beta)}{2}
\int_{\mathbb{R}^N}(I_\beta\ast|u_1^n|^{r_1})|u_2^n|^{r_2}dx \\
&\to  \sum_{i=1}^2\frac{\mu_i}{2p_i}(Np_i-N-\alpha)\int_{\mathbb{R}^N}(I_\alpha
 \ast|u_i^n|^{p_i})|u_i^n|^{p_i}dx+\Theta\frac{Nr-2(N+\beta)}{2}
 \int_{\mathbb{R}^N}(I_\beta\ast|u_1|^{r_1})|u_2|^{r_2}dx.
\end{align*}
Combining $P(u_1,u_2)=0$, we have
\[
\lim_{n\to +\infty}a\sum_{i=1}^2\| \nabla u_i^n\|_2^2
 +b\sum_{i=1}^2\| \nabla u_i^n\|_2^4
= a\sum_{i=1}^2\| \nabla u_i\|_2^2+b\sum_{i=1}^2\| \nabla u_i\|_2^4.
\]
Hence, $F(u_1^n,u_2^n)=F(u_1,u_2)$, and then $F(u_1,u_2)=\delta(d_1,d_2)$.
\end{proof}

\begin{proof}[Proof of Theorem \ref{thm1.3}]
By Lemma \ref{lmk}, we only need to prove $(u_1,u_2)\in S_r(d_1)\times S_r(d_2)$. 
Since $(u_1,u_2)$ is solution of \eqref{yy}, by Lemma \ref{zjngz}, we have 
$\lambda_1<0$ or $\lambda_2<0$. Without loss of generality, we can assume 
$\lambda_1<0$. By Lemma \ref{yuk9},
we obtain that $u_1^n\to  u_1$ in $H_r^1(\mathbb{R}^N)$ and then $u_1\in S_r(d_1)$. 
Suppose by contradiction that $\lambda_2\geq0$, then
\[
 -\Delta u_2=\lambda_2 u_2+\mu_2(I_\alpha\ast|u_2|^{p_2})|u_2|^{p_2-2}u_2+
 \Theta r_2(I_\beta\ast|u_1|^{p_1})|u_2|^{p_2-2}u_2\geq0.
\]
By Lemma \ref{Liuv}, we have $u_2=0$. Therefore, $F(u_1,u_2)=F(u_1,0)$, and
$u_1\in S_r(d_1)$ satisfies \eqref{sig}.
By Lemma \ref{byg}, $F(u_1,0)=\sigma_{p_1}^{\mu_1}(d_1)<0$ if $a=1$,
$b=0$, $N=3$ and $p_1=p_2=r_1=r_2=\alpha=\beta=2$, which is in contradiction to the
fact $F(u_1,0)=\delta(d_1,d_2)>0$. Thus, $\lambda_2<0$, and then $u_2\in S_r(d_2)$.
By the maximum principle, we obtain the $u_1,u_2>0$ in $\mathbb{R}^3$.
\end{proof}

\subsection*{Acknowledgements}
This study was funded by the Fujian Provincial Natural Science
of China (No. 2025J08201).
Zhewen Chen thanks Professor Jianqing Chen, at the Fujian normal University, 
for his valuable comments in preparing this article.

\begin{thebibliography}{10}

\bibitem{Kf1} A. Arosio, S. Panizzi;
\emph{On the well-posedness of the Kirchhoff string}. 
Trans. Am. Math. Soc., 348 (1996), 305-330.

\bibitem{Huij4} T. Bartsch, H. W. Li, W. M. Zou;
\emph{Existence and asymptotic behavior of normalized ground states for Sobolev
critical Schr\"{o}dinger systems}, Calc. Var. Partial Differential Equations, 62 (2023) 9.


\bibitem{Huij1} T. Bartsch, L. Jeanjean, N. Soave;
\emph{Normalized solutions for a system of coupled cubic Schr\"{o}dinger equations
on $\mathbb{R}^3$}, J. Math. Pures Appl., 106 (2016) 583-614.

\bibitem{Huij5} T. Bartsch, N. Soave;
\emph{A nature constraint approach to normalized solutions of nonlinear 
Schr\"{o}dinger equations and systems}, J. Funct. Anal., 272 (2017) 4998-5037.

\bibitem{Huij2} T. Bartsch, N. Soave;
\emph{Multiple normalized solutions for a competing system of Schr\"{o}dinger equations},
 Calc. Var. Partial Differential Equations 58 (2019) 22.

\bibitem{Huij3} T. Bartsch, X. X. Zhong, W. M. Zou;
\emph{Normalized solutions for a coupled Schr\"{o}dinger system}, Math. Ann.,
380 (2021) 1713-1740.

\bibitem{Arm2} H. Berestycki, P. L. Lions;
\emph{Nonlinear scalar field equations II},
 Arch. Ration. Mech. Anal., 82(1983) 347-376.

\bibitem{kf2} X. Cao, J. Xu, J. Wang;
\emph{The existence of solutions with prescribed $L^2$-norm for Kirchhoff type system}. 
J. Math. Phys., 58 (2017) 041502.

\bibitem{BZd27} I. E. Dzyaloshinskii, E. M. Lifshitz, Lev P. Pitaevskii;
\emph{General theory of van der Waals' forces}. Sov. Phys. Usp., 4 (1961) 153.

\bibitem{JFx} Q. P. Geng, Y. Y. Tu, J. Wang;
\emph{Existence and multiplicity of the positive normalized solutions to the
coupled Hartree-Fock type nonlocal elliptic system}, 
J. Fixed Point Theory Appl., 24 (2022) 83.

\bibitem{banqun} M. Ghimenti, J. Van Schaftingen;
\emph{Nodal solutions for the Choquard equation}, J. Funct. Anal. 271
(2016) 107-135.

\bibitem{sben} N. Ghoussoub;
\emph{Duality and perturbation methods in critical point theory}, Cambridge Tracts
in Mathematics, vol. 107 (Cambridge University Press, 1993).

\bibitem{Huij6} T. X. Gou, L. Jenajean;
\emph{Multiple positive normalized solutions for nonlinear Schr\"{o}dinger systems},
Nonlinearity, 31 (2018) 2319-2345.

\bibitem{guo1} Z. Y. Guo, L. J. Zhao;
\emph{Existence of ground states for fractional Choquard-Kirchhoff
equations with magnetic fields and critical exponents}, 
Period. Math. Hungar. 87 (2023) 468-483.

\bibitem{guo2} Z. Y. Guo, T. Q. Zhang;
\emph{Normalized solutions of fractional Kirchhoff equations in the
defocusing case}, Electronic J. Differ. Equ., 2025 (2025) No. 75, 1-16.

\bibitem{Kf2} J. Q. Hu, A. M. Mao;
\emph{Normalized solutions to nonlocal Schr\"{o}dinger systems with $L^2$-subcritical
 and supercritical nonlinearities}, J. Fixed Point Theory Appl., 25 (2023), 77.

\bibitem{Kf5} Q. H. He, Z. Y. Lv, Z. W. Tang;
\emph{The existence of normalized solutions to the Kirchhoff equation with potential 
and Sobolev critical nonlinearities}, J. Geom. Anal., 33 (2023), 236.

\bibitem{Kf6} Q. H. He, Z. Y. Lv, Y. M. Zhang, X. X. Zhong;
\emph{Existence and blow up behavior of positive normalized solution to the Kirchhoff
 equation with general nonlinearities:  Mass super-critical case},
  J. Differ. Equ., 356 (2023), 375-406.

\bibitem{em} N. Ikoma;
\emph{Compactness of minimizing sequences in nonlinear Schr\"{o}dinger systems under
multiconstraint conditions}, Adv. Nonlinear Stud., 14 (2014), 115-136.

\bibitem{97Na} L. Jeanjean;
\emph{Existence of solutions with prescribed norm for semilinear elliptic equations}.
Nonlin. Analysis, 28 (1997), 1633-1659.

\bibitem{Kf4} G. Kirchhoff;
\emph{Mechanik}. Teubner, Leipzig, 1883.

\bibitem{lieb} E. H. Lieb, M. Loss;
\emph{Analysis}, Gradute Studies in Mathematics, AMS, Providence, Rhode island,
2001.

\bibitem{Li4} J. Lions;
\emph{On some questions in boundary value problems of mathematical physics}. 
North-Holland Math. Stud., 30 (1978), 284-346.

\bibitem{ygg} Z. Liu;
\emph{Multiple normalized solutions for Choquard equations involving Kirchhoff type 
perturbation}, Topol. Methods Nonlinear Anal., 54 (2019) 297-319.

\bibitem{Wyu} L. Ma, L. Zhao;
\emph{Classiflcation of positive solitary solutions of the nonlinear Choquard equation},
 Arch. Ration. Mech. Anal., 195 (2010), 455-467.

\bibitem{Mza} V. Moroz, J. Van Schaftingen;
\emph{Existence of groundstates for a class of nonlinear Choquard
equations}, Trans. Amer. Maths. Soc., 367 (2015), 6557-6579.

\bibitem{cfrv} S. G. Porsev, A. Derevianko;
\emph{High-accuracy calculations of dipole, quadrupole, and octupole electric dynamic 
polarizabilities and van der Waals coefficients $C_6$, $C_8$, and $C_{10}$ 
for alkaline-earth dimers}, J. Exp. Theor. Phys. 102 (2006), 195-205.

\bibitem{GN23} M. I. Weinstein;
\emph{Nonlinear Schr\"{o}dinger equations and sharp interpolation estimates}.
Comm. Math. Phys. 87 (1983), 567-576.

\bibitem{Kf3} Z. Yang;
\emph{Normalized ground state solutions for Kirchhoff type systems}. 
J. Math. Phys., 62 (2021), 031504.

\bibitem{VWp3} Y. Zheng, A. Narayanaswamy;
\emph{Lifshitz theory of van der Waals pressure in dissipative media}. 
Phys. Rev. A, 83 (2011), 042504.

\bibitem{linjie} H. Zhang, J. J. Zhang, X. X. Zhong;
\emph{Normalized solutions for critical Choquard systems}, arXiv:2307.01483v2

\end{thebibliography}
\end{document}

